\documentclass[12pt,a4paper,oneside]{memoir}

% ===============================
% 1. PACKAGES
% ===============================
\usepackage{times}
\usepackage[T1]{fontenc}
\usepackage{xcolor}
\usepackage{graphicx}
\usepackage{amsmath,amssymb}
\usepackage{booktabs}
\usepackage{hyperref}
\usepackage{caption}
\usepackage[backend=biber,style=ieee]{biblatex}
\usepackage{siunitx}
\usepackage{array}
\usepackage{longtable}
\usepackage{placeins}
\usepackage{pdflscape}
\usepackage{lipsum}        % dummy text
\usepackage[linesnumbered,ruled,vlined]{algorithm2e}
\usepackage{float}
%
\usepackage{listings}
% --- Custom Colors for Python Beautification ---
\definecolor{codegreen}{rgb}{0,0.6,0}
\definecolor{codegray}{rgb}{0.5,0.5,0.5}
\definecolor{codepurple}{rgb}{0.58,0,0.82}
\definecolor{backcolour}{rgb}{0.96,0.96,0.97}
\definecolor{codeblue}{rgb}{0.1,0.2,0.8}
% --- Global Listings Style ---
\lstset{
	backgroundcolor=\color{backcolour},   
	commentstyle=\color{codegreen},
	keywordstyle=\color{codeblue}\bfseries,
	numberstyle=\tiny\color{codegray},
	stringstyle=\color{codepurple},
	basicstyle=\ttfamily\scriptsize,
	breakatwhitespace=false,         
	breaklines=true,                 
	captionpos=b,                    
	keepspaces=true,                 
	numbers=left,                    
	numbersep=8pt,                  
	showspaces=false,                
	showstringspaces=false,
	showtabs=false,                  
	tabsize=1,
	frame=single,
	rulecolor=\color{codegray!30},
	language=Python
}

% ===============================
% 2. BIBLIOGRAPHY SETUP
% ===============================
\addbibresource{References.bib}

% ===============================
% 3. COLORS & SPACING
% ===============================
\definecolor{deepred}{RGB}{152,50,34}
\linespread{1.25}
\setlength{\parindent}{0pt}
\setlength{\parskip}{6pt}
\hyphenpenalty=10000
\exhyphenpenalty=10000
\emergencystretch=3em

% ===============================
% 4. PAGE LAYOUT
% ===============================
\setlrmarginsandblock{1.1in}{1.1in}{*}
\setulmarginsandblock{1in}{1in}{*}
\checkandfixthelayout
\pagestyle{plain}

% ===============================
% 5. NUMBERING & TOC DEPTH
% ===============================
\setsecnumdepth{subsubsection}
\maxsecnumdepth{subsubsection}
\settocdepth{subsection}
\setcounter{tocdepth}{2}

% ===============================
% 6. CHAPTER STYLE (Centered & ALL CAPS)
% ===============================
\makechapterstyle{sdpstyle}{
	\renewcommand{\printchaptername}{}
	\renewcommand{\chapternamenum}{}
	\renewcommand{\printchapternum}{\centering\large\bfseries \thechapter.\space}
	\renewcommand{\chaptitlefont}{\large\bfseries\centering}
	\renewcommand{\afterchapternum}{}
	\renewcommand{\printchaptertitle}[1]{%
		\chaptitlefont\MakeUppercase{##1}%
	}
}
\chapterstyle{sdpstyle}

% ===============================
% 7. HEADING STYLES
% ===============================
\hypersetup{
	colorlinks=false,
	linkcolor=black,
	citecolor=black,
	filecolor=black,
	urlcolor=black,
	linktoc=all,
	pdfborder={0 0 0},
}
\setsecheadstyle{\color{black}\normalsize\bfseries\scshape}
\setsubsecheadstyle{\color{black}\normalsize\bfseries\scshape}
\setsubsubsecheadstyle{\color{black}\bfseries\normalshape}
\renewcommand{\thesubsubsection}{\textcolor{black}{\thesubsection.\arabic{subsubsection}}}

% --- Paragraph heading style (memoir native) ---
\setparaheadstyle{\bfseries}
\setafterparaskip{1.5ex plus .4ex minus .2ex}
\setbeforeparaskip{3.25ex plus 1ex minus .2ex}

% ===============================
% 8. TOC STYLING
% ===============================
\renewcommand{\printtoctitle}[1]{\centering\large\bfseries \MakeUppercase{#1}}
\renewcommand{\cftchapterfont}{\scshape\bfseries\color{black}}
\renewcommand{\cftchapterpagefont}{\color{black}}
\renewcommand{\cftsectionfont}{\scshape\color{black}}
\renewcommand{\cftsectionpagefont}{\color{black}}
\renewcommand{\cftsubsectionfont}{\normalshape\color{black}}
\renewcommand{\cftsubsectionpagefont}{\color{black}}

% ===============================
% 9. UNNUMBERED RED HEADING COMMAND
% ===============================
\newcommand{\redheading}[1]{%
	\vspace{6pt}
	{\color{black}\bfseries\normalshape #1}\par
	\noindent
	\vspace{0pt}%
}

% ===============================
% 10. FIGURE & TABLE CAPTIONS
% ===============================
\captionnamefont{\scshape}
\captiontitlefont{\upshape}
\captiondelim{.\ }
\setlength{\belowcaptionskip}{10pt}
\setlength{\abovecaptionskip}{10pt}

% ===============================
% 11. CONTENTS NAME
% ===============================
\renewcommand{\contentsname}{CONTENTS}

% ===============================
% 12. ALGORITHMS (algorithm2e)
% ===============================
\renewcommand{\listalgorithmcfname}{List of Algorithms}

% ===============================
% 13. PAGE NUMBERING & LIST STYLING
% ===============================
\frontmatter
\pagenumbering{roman}

% --- Figure Spacing (Tightened) ---
\renewcommand{\cftfigurename}{Figure~}
\renewcommand{\cftfigureaftersnum}{.}
\setlength{\cftfigurenumwidth}{3.0em} 
\setlength{\cftfigureindent}{0em}    

% --- Table Spacing (Tightened) ---
\renewcommand{\cfttablename}{Table~}
\renewcommand{\cfttableaftersnum}{.}
\setlength{\cfttablenumwidth}{3.0em} 
\setlength{\cfttableindent}{0em}    

% --- Algorithm List Styling ---
\makeatletter
% Force "Algorithm" prefix into the list
\let\oldl@algocf\l@algocf
\renewcommand{\l@algocf}[2]{\oldl@algocf{Algorithm #1}{#2}}

% Direct memoir-style length setting for Algorithms
\AtBeginDocument{
	\newlength{\cftalgocfnumwidth}
	\setlength{\cftalgocfnumwidth}{6em}
	\newlength{\cftalgocfindent}
	\setlength{\cftalgocfindent}{0em}
}
\makeatother

% ===============================
% 14. APPENDIX SETUP 
% ===============================
\makeatletter
\g@addto@macro\appendix{%
	% 1. Force A, B, C numbering
	\renewcommand{\thechapter}{\Alph{chapter}}
	
	% 2. Fix the Header on the actual page
	\renewcommand{\printchapternum}{%
		\centering\large\bfseries APPENDIX \thechapter.\space%
	}
	
	% 3. Override the chapter command
	\def\chapter{\@ifstar{\@schapter}{\@dblarg\@manualappendixchapter}}
}

\def\@manualappendixchapter[#1]#2{%
	\cleardoublepage % <--- THIS FORCES THE NEW PAGE
	\refstepcounter{chapter}%
	% Writes "Appendix A. Title" to the TOC
	\addcontentsline{toc}{chapter}{\scshape Appendix \thechapter.\hspace{1em}#1}%
	\chaptermark{#1}%
	\if@twocolumn
	\@topnewpage[\@makechapterhead{#2}]%
	\else
	\@makechapterhead{#2}%
	\@afterheading
	\fi
}
\makeatother
%
% ==================================================
% DOCUMENT START
% ==================================================
\begin{document}

% --- TITLE PAGE ---
\begin{titlingpage}
	\centering
	\linespread{1.0}
	\vspace*{0.5in}
	{\large\textbf{\MakeUppercase{Design and Implementation of a Scalable IoT-Based Acoustic Detection System for Early Red Palm Weevil Infestation}}\par}
	\vspace{0.7in}
	{\large\textbf{\textsc{CMP/COE490: Senior Design Project}}\par}
	\vspace{0.4in}
	Submitted in Partial Fulfillment of the\\
	Requirements for the Degree of\\
	\vspace{0.4in}
	\textbf{\textsc{Bachelor of Science in Computer Science and Engineering}}
	\vspace{0.4in}

	by

	\vspace{0.2in}

	\begin{tabular}{p{2.2in} r}
		Khalid Hamdan         & 99009  \\
		Abdurrahman Maricar   & 96961  \\
		Kareem Mahmoud        & 97349  \\
		Haitham Ayyash        & 103430 \\
	\end{tabular}

	\vspace{0.4in}

	\textit{Advisors}\\[0.15in]
	Dr.\ Mohamed Alhajri\\
	Dr.\ Taha Landolsi

	\vfill

	\makebox[\textwidth][c]{%
		\includegraphics[width=0.175\textwidth]{AUSlogo.png}
	}

	\vspace{0.1in}

	Department of Computer Science and Engineering\\
	American University of Sharjah\\[0.15in]
	Spring 2026
\end{titlingpage}

% --- ACKNOWLEDGMENTS ---
\chapter*{Acknowledgments}
\thispagestyle{plain}
\addcontentsline{toc}{chapter}{Acknowledgments}
\vspace*{0.4in}
\begin{center}
\itshape
With sincere gratitude
\end{center}
\vspace{0.2in}
\begin{center}
On behalf of our Senior Design team, we would like to extend our sincere thanks to our advisor, Dr.\ Mohamed Alhajri, and our co-advisor, Dr.\ Taha Landolsi, for their guidance, dedication, and continuous support throughout this project.

We are also grateful to the American University of Sharjah and the College of Engineering for providing us with the opportunity and support to pursue this work.

Last but not least, we would like to express our heartfelt thanks to our families for their constant encouragement, patience, and unwavering support.
\end{center}

\vspace{0.5in}
\begin{flushright}
\textit{Respectfully submitted,}\\[0.1in]
\textbf{The Senior Design Team}
\end{flushright}
\clearpage

% --- ABSTRACT ---

\chapter*{Abstract}
\thispagestyle{plain}
\addcontentsline{toc}{chapter}{Abstract}
Palms face significant threats from the Red Palm Weevil (RPW), making RPW management challenging due to the internal development of larvae within the trunk, which often go unnoticed until infestations escalate. This paper presents an innovative, scalable Internet of Things (IoT) system for the acoustic detection and early identification of RPW. The proposed system is designed to use low-cost, trunk-mounted piezoelectric sensors to capture vibrations caused by larval activity within the palm. These sensors are based on ESP32 microcontrollers and LoRa communication modules and are intended to form a wireless sensor network that transmits data packets to a central hub for cloud storage, analysis, and alert generation.

To ensure reliable signal detection, the proposed design incorporates acoustic preprocessing techniques, including bandpass filtering, SNR gating, and MFCC extraction, followed by a deep learning-based inference pipeline intended to distinguish infestation-related acoustic signals from environmental noise. A web dashboard is planned for real-time monitoring of tree health, node activity, and potential infestation events. The architecture is designed to be cost-effective and modular, with the goal of supporting large-scale agricultural deployment. Additionally, the proposed system aims to facilitate early interventions, minimize pesticide usage, and alleviate economic losses stemming from delayed inspection processes, thus outlining a viable precision agriculture approach for non-invasive, real-time monitoring of palm tree health.

\clearpage

% --- CONTENTS (The TOC itself) ---
\tableofcontents* \clearpage

% --- LIST OF FIGURES ---
\listoffigures
\clearpage

% --- LIST OF TABLES ---
\listoftables
\clearpage

% --- GLOSSARY ---

\chapter*{Glossary}
\thispagestyle{plain}
\addcontentsline{toc}{chapter}{Glossary}
\renewcommand{\arraystretch}{1.3}
\begin{tabular}{@{}p{0.12\linewidth} p{0.85\linewidth}@{}}
	\textbf{ADC} & Analog-to-Digital Converter \\
	\textbf{AI} & Artificial Intelligence \\
	\textbf{ANN} & Artificial Neural Network \\
	\textbf{API} & Application Programming Interface \\
	\textbf{Bi-LSTM} & Bidirectional Long Short-Term Memory \\
	\textbf{CNN} & Convolutional Neural Network \\
	\textbf{DAS} & Distributed Acoustic Sensing \\
	\textbf{ESP32} & Espressif 32-bit Microcontroller \\
	\textbf{IoT} & Internet of Things \\
	\textbf{IPM} & Integrated Pest Management \\
	\textbf{JSON} & Java Script Object Notation \\
	\textbf{JWT} & JSON Web Token \\
	\textbf{LoRa} & Long Range Radio \\
	\textbf{ML} & Machine Learning \\
	\textbf{RPW} & Red Palm Weevil \\
	\textbf{SNR} & Signal-to-Noise Ratio \\
	\textbf{UI} & User Interface \\
	\textbf{VOC} & Volatile Organic Compound \\
	\textbf{Wi-Fi} & Wireless Fidelity \\
	\textbf{WSN} & Wireless Sensor Network \\
	\textbf{$\Phi$-OTDR} & Phase-Sensitive Optical Time-Domain Reflectometry \\
\end{tabular}

\clearpage

\mainmatter
\pagenumbering{arabic}

\chapter{Introduction}
\label{ch:Introduction}
 
% ===============================
% CHAPTER 1: INTRODUCTION
% ===============================
 
\section{Motivation}
\label{sec:Motivation}
 
The driving force behind this project is the increasing effect of Red Palm Weevil (RPW) infestations on agriculture and the economy. This section assesses the current state of the RPW infestation and evaluates the existing detection technologies.
 
\subsection{Current Situation of Red Palm Weevil Infestation}
 
The RPW (Rhynchophorus ferrugineus) is currently recognized as one of the most destructive and invasive pests threatening palm species worldwide, especially in the Middle East, North Africa, and the Mediterranean basin. The economic losses associated with severely infested palms in Gulf countries have been estimated in the millions of U.S. dollars annually \cite{saleem2022,alaufi2026}. The primary challenge lies in the ``hidden'' nature of the pest. For most of the larval life cycle, it remains concealed inside the trunk while feeding, making visual inspection unreliable until the damage is already advanced \cite{saleem2022,alaufi2026}. Visible symptoms such as holes in the trunk or crown, oozing brown viscous liquid, and collapsed crowns usually indicate a severe infestation that has progressed to the point where the palm tree is already dying and difficult or impossible to save \cite{saleem2022,alaufi2026}. These limitations motivate the development of a cost-effective, reliable, and non-invasive early detection system capable of deployment across large plantations.
 
\begin{figure}[H]
\centering
\includegraphics[width=0.75\textwidth]{figures/fig1_rpw_lifecycle.jpeg}
\caption{Schematic Diagram Depicting the Life Cycle of RPW \cite{alaufi2026}}
\label{fig:rpw_lifecycle}
\end{figure}
 
\begin{figure}[H]
\centering
\includegraphics[width=0.75\textwidth]{figures/fig2_viscous_fluid.jpeg}
\caption{Viscous fluids ooze from the tunnels due to larvae inside the palm tree \cite{alaufi2026}}
\label{fig:viscous_fluid}
\end{figure}

\FloatBarrier
 
\subsection{Existing Solutions and Their Limitations}
 
Several methods have been explored to tackle the early detection of the RPW infestations, yet most existing approaches suffer from limitations related to scalability, invasiveness, cost, or detection accuracy. Visual and manual inspection remain the most adopted techniques; however, it is labor intensive, expensive, and subjective to the skills of the laborer. Even trained experts can inspect only a limited number of trees per day, and this method is particularly ineffective for detecting early-stage infestations due to the concealed larval activity inside the trunk \cite{saleem2022,alaufi2026}.
 
\begin{figure}[H]
\centering
\includegraphics[width=0.75\textwidth]{figures/fig3_infected_crown.jpeg}
\caption{The crown of an infested young date palm (left) ejects chewed-up fibers. Internal palm tissue
shattered because of RPW Larvae consumption \cite{alaufi2026}}
\label{fig:infested_crown}
\end{figure}
 
Another standard method is to use pheromone trapping, which provides valuable information on adult weevil population trends and overall infestation pressure \cite{shukla2018,na2024}. However, traps only work on adult insects and yield no information about actual larval activity occurring inside the palm trees \cite{saleem2022}. Likewise, thermal imaging and remote sensing methods attempt to measure physiological stress signals or internal heat signatures, but these systems are frequently influenced by environmental conditions and generally provide lower reliability for large-scale field assessment than acoustic methods \cite{saleem2022,alaufi2026}. Conventional acoustic sensors are therefore a more promising early-warning approach \cite{karar2021,karar2022}. Piezoelectric probe-based devices are highly sensitive to larval gnawing sounds and, under controlled conditions, may detect infestations within the first two weeks of larval development \cite{ashry2020,ashry2022cnn}. Although such systems are typically invasive (e.g., they involve drilling the tree trunk) \cite{pinhas2008,karar2022}, they provide reliable and well-established detection performance. Despite their higher installation and maintenance costs, they remain suitable for targeted monitoring in semi-controlled environments \cite{karar2022}.

\begin{figure}[H]
\centering
\includegraphics[width=0.5\textwidth]{figures/fig4_treevibes.jpeg}
\caption{TreeVibes sensing device for acquiring sounds of RPW larvae \cite{rigakis2021}}
\label{fig:treevibes}
\end{figure}
 
On the other hand, the combination of Distributed Acoustic Sensing (DAS) and Convolutional Neural Networks (CNNs) is the path with higher potential to meet project goals \cite{ashry2020,wang2021,ashry2022cnn}. DAS technology turns a conventional optical fiber into an external vibration sensor, allowing a single cable wrapped non-invasively around trunks to monitor many trees concurrently \cite{ashry2020,mao2021}. In conjunction with CNN-based classifiers, this method has reported classification accuracies between about 97\% and 99.9\% under controlled and semi-controlled studies \cite{wang2021,ashry2022cnn}. However, due to the high cost for installation, the focus will be on the implementation and evaluation of the conventional acoustic sensing approach as a practical and experimental solution.
 
 
 
\section{Project Description}
\label{sec:Description}
 
While DAS represents the highest-performing solution in the literature, its installation cost renders it impractical for large-scale deployment. This project aims to develop an affordable, IoT-based acoustic detection system for larvae (of RPW) in an early infestation stage. Instead of expensive distributed fiber-optic technology, the new system will be made from arrays of low-cost piezoelectric contact sensors mechanically coupled to palm trunks, capable of detecting internal vibrations produced by larval feeding and movement. By combining low-cost sensor devices, deep learning classifiers and cloud-based data analytics, this work intends to investigate a scalable and affordable open solution for real-time plantation monitoring at large-scale agroecological landscapes. The architecture follows a three-tier design: trunk-mounted piezoelectric sensing nodes transmit compact ~214-byte feature packets consisting of quantized MFCC features, burst statistics, and metadata via LoRa to a cluster head, which aggregates and forwards data to a cloud backend for deep learning classification. The proposed system will be able to detect an RPW infestation earlier than visual inspection, before severe structural damage occurs \cite{ashry2020,ashry2022cnn,karar2022}. The system is designed around low-cost components, a modular IoT infrastructure, and automated analysis tools, to be low-cost and widely deployed. Thus, continuous remote monitoring could help reduce the considerable amount of field labor involved in pest control, and supply information to help reduce pesticide use, both of which would be helpful economically and environmentally \cite{saleem2022,karar2022}. As a secondary contribution, this project evaluates the recently introduced Mamba selective state-space model architecture for RPW acoustic classification, comparing it against established CNN and CNN-LSTM baselines.
 
\section{Problem Statement}
\label{sec:Statement}
 
Red Palm Weevil (RPW) is one of the most destructive pests affecting palm species, particularly date palms in the Middle East and North Africa Region. The larvae develop and feed internally within the trunk, causing progressive structural damage that remains undetectable through visual inspection during early infestation stages. By the time external symptoms such as trunk perforations, oozing sap, or crown collapse become visible, the palm tree is often severely damaged and difficult to recover. This concealed life cycle creates a critical gap between infestation onset and detection, resulting in delayed intervention, increased tree mortality, economic losses, and excessive pesticide use. This project aims to design, implement, and validate a low-cost, IoT-based acoustic detection system for early identification of RPW larvae inside palm trunks. The primary objective is to enable real-time monitoring of internal trunk vibrations generated by larval feeding and movement, thereby detecting infestations prior to visible external symptoms.
 
The main solution proposed in this project is the development of a distributed network of trunk-mounted piezoelectric acoustic sensors integrated with wireless communication modules and cloud-based deep learning classification framework. Acoustic signals captured from the trunk are processed into spectrogram representation and analyzed using convolutional neural network architectures to distinguish larval impulsive sounds from environmental noise. By combining low-cost sensing hardware, wireless data transmission, and automated intelligent classification, the system seeks to provide a scalable, non-destructive, and economically viable early detection solution for large-scale date palm plantations. In the following section, the challenges identified here are translated into clear, measurable design objectives.
 
\section{Design Objectives}
\label{sec:Objectives}
 
The primary goal of this project is to design, implement, and validate a low-cost, IoT-based acoustic sensing system for the early detection of Red Palm Weevil larvae in palm trees. The system integrates piezoelectric trunk-mounted sensors with deep learning-based signal classification to provide non-invasive, real-time, and scalable monitoring. The following objectives define the specific, tangible, and measurable deliverables of the project.
 
\subsection{Development of a Low-Cost Acoustic Sensing Unit}
 
The project shall design and prototype a trunk-mounted acoustic sensing unit capable of detecting internal larval feeding vibrations within a radius of 1.5 meters to 2.0 meters from the sensor location. The sensing module shall incorporate a high-sensitivity piezoelectric element coupled to a minimally invasive metallic waveguide to enhance vibration transmission. The total hardware cost per sensing node shall not exceed a predefined affordability threshold suitable for large-scale plantation deployment. Each sensing unit shall include a microcontroller platform capable of sampling acoustic signals at an adequate frequency of resolution for spectrogram generation. The unit shall be integrated with a LoRa wireless communication module to enable data forwarding to a cluster head, which then relays data to a centralized cloud server over Wi-Fi or cellular backhaul. A functional prototype of the sensing node shall be completed and validated under laboratory conditions before field deployment.
 
\subsection{Implementation of Wireless Sensor Network Architecture}
 
The system shall establish a scalable Wireless Sensor Network architecture that supports simultaneous operation of multiple sensing nodes across a plantation. The network shall enable reliable transmission of acoustic data packets to a centralized cloud-based server with minimal data loss.
 
Power management shall be optimized to allow extended battery-powered operation using low-power electronics and duty-cycled deep sleep, with optional solar charging for long-term deployments. Each node shall be enclosed in a rugged outdoor casing to ensure environmental protection against heat, humidity, and dust. System performance shall be validated in terms of communication stability, node synchronization, and operational uptime under simulated field conditions.
 
 
\subsection{Development of a Deep Learning-Based Acoustic Classifier}
 
The project shall design and implement a deep learning model capable of differentiating larval impulsive feeding sounds from environmental noise sources such as wind, birds, and agricultural machinery. The deep learning classifier architecture shall be trained and validated using labeled acoustic datasets. Because missed infestations directly undermine the early-detection objective of this project, model selection shall prioritize recall, with the false negative rate and F1 score reported as supporting metrics. Model robustness shall be evaluated against varying noise conditions to ensure reliable operation in real-world plantation environments. To further validate the model's robustness, 10-fold cross-validation will be performed, with variance across the trained models evaluated to ensure it is less than 3\% in the chosen primary metric.
 
 
\subsection{Integration of Cloud-Based Data Processing and Monitoring Platform}
 
A centralized cloud infrastructure shall be developed to store, process, and analyze incoming MFCC feature streams transmitted by sensing nodes. The platform shall support automated validation, normalization, model inference, and alert generation when larval activity is detected. A user interface dashboard shall be implemented to provide plantation managers with real-time monitoring capabilities, including visualization of sensor status, detection events, and geographic mapping of suspected infestation zones. The system shall generate automated alerts through predefined communication channels when detection confidence exceeds a specified threshold.
 
\subsection{Validation Through Field Testing}
 
The complete system, including hardware, communication infrastructure, and classification software, shall undergo field validation in a controlled plantation environment. Performance evaluation shall include detection accuracy, false positive rate, latency between larval activity and alert generation, and system reliability over continuous monitoring periods. Comparative analysis shall be conducted against conventional inspection practices to demonstrate earlier detection capability prior to visible external symptoms. Results shall be documented to verify that the system satisfies its functional objectives and meets the intended scalability and affordability targets
 

\subsection{Deployment and Operational Integration}
 
The project shall define procedures for large-scale deployment of the acoustic detection system across commercial plantations. Deployment guidelines shall include standardized installation protocols for mounting sensor nodes on palm trunks, ensuring proper acoustic coupling, and verifying wireless connectivity. The system shall be configured to operate as an integrated network, with each node transmitting data reliably to the centralized cloud server. Deployment objectives shall also include site-specific calibration of sensors, validation of signal quality after installation, and operational verification under actual environmental conditions, including wind, rain, and varying temperatures. Training materials and operational manuals shall be developed to enable plantation staff to perform basic maintenance, node replacement, and troubleshooting. The deployment strategy shall ensure scalability, modularity, and rapid expansion capability, enabling full plantation coverage with minimal labor while maintaining system accuracy and real-time monitoring capabilities.
 
\section{Design Constraints}
\label{sec:Constraints}
 
In the context of deploying an acoustic microphone-based detection system for early Red Palm Weevil (RPW) detection, the following design constraints have been identified based on the technical limitations documented in the literature and the practical realities of real-world farm deployment.
 
\subsection{Environmental Noise Limitations}
 
A fundamental constraint of this project is the degradation in system performance under noisy outdoor conditions. Wind, rain, and ambient farm sounds such as bird activity and machinery generate acoustic interference that overlaps with the feeding frequency range of RPW larvae (100--2500 Hz), leading to a significant increase in false alarms. Contact microphone and piezoelectric probe systems are particularly susceptible to this interference when deployed in open-air environments, and no universal noise cancellation solution has been demonstrated to fully eliminate this limitation under all field conditions \cite{pinhas2008,karar2022,ashry2022cnn}.
 
\subsection{Restriction to Early-Stage Larvae Detection}
 
The acoustic microphone system is designed to detect RPW during the larval stage, specifically when larvae are actively feeding and producing chewing sounds inside the tree trunk. Adult weevils and eggs do not produce comparable acoustic emissions, which means the detection window is restricted to a narrow biological stage. Moreover, larvae do not produce feeding sounds continuously; intermittent feeding behavior reduces the reliability of single-measurement detection windows and may require extended monitoring periods to accumulate sufficient acoustic evidence of infestation \cite{sutanto2023,wang2021}.
 
\subsection{Invasiveness and Tree Integrity Constraints}
 
A significant practical constraint of probe-based acoustic systems is their invasive nature. Deploying piezoelectric contact microphones or waveguide rods typically requires drilling a small hole into the tree trunk to improve acoustic coupling between the sensor and the internal wood structure. This procedure risks introducing entry points for secondary pests, pathogens, and fungal infections, potentially causing additional harm to the trees being monitored. Repeated insertions across large plantations compound this risk and may be objectionable from an agronomic or regulatory standpoint \cite{pinhas2008,karar2022}.
 
\subsection{Sensor Coverage and Spatial Limitation}
 
The effective detection radius of a single piezoelectric or contact microphone probe is typically limited to approximately 1.5 to 2.0 meters from the insertion point \cite{pinhas2008,karar2022}. This means that taller trees or trees with large trunk diameters may require multiple sensors placed at different heights to ensure adequate coverage of the entire trunk volume where larvae may be present. Scaling this approach to large farms with hundreds or thousands of trees requires a proportional increase in the number of sensors, cabling, and wireless communication modules, significantly increasing infrastructure costs and deployment complexity \cite{pinhas2008,karar2022,saleem2022}.
 
\subsection{Computational and Real-Time Processing Constraints}
 
The machine learning models used for classification including Convolutional Neural Networks (CNNs), Bidirectional Long Short-Term Memory (Bi-LSTM) networks, and statistical models such as Gaussian Mixture Models (GMMs) require pre-trained weights and dedicated processing hardware. Feature extraction pipelines involving Mel-Frequency Cepstral Coefficients (MFCCs), Short-Time Fourier Transform (STFT) spectrograms, or Linear Predictive Coding (LPC) coefficients impose additional computational overhead. While classification accuracies exceeding 97\% have been reported under controlled laboratory conditions, extending these models to diverse real-world noise environments requires retraining on larger and more representative datasets, which may not be readily available \cite{boulila2023,karar2022,saranya2025}.
 
\subsection{Cost and Scalability Constraints}
 
Although individual piezoelectric sensors and microphone modules are relatively inexpensive compared to imaging technologies such as X-ray tomography, the requirement to assign at least one and often multiple sensors per tree, combined with the need for wireless communication nodes, power supplies, and data acquisition infrastructure, results in a total system cost that scales unfavorably with farm size. For large commercial plantations with thousands of trees, the cumulative hardware, installation, and maintenance costs may become prohibitive without institutional support or subsidy programs \cite{karar2022,saleem2022}.
 
\section{Project Impact}
\label{sec:Impact}
 
\subsection{Economic Impact}
 
The deployment of an IoT-based networked acoustic system for early RPW detection carries substantial economic implications across multiple levels of the agricultural value chain, from individual farmers to national economies. At the macro level, RPW infestation generates annual losses of millions of US dollars worldwide through destroyed palm crops and pest control expenditures \cite{ashry2020,saleem2022}. In Gulf countries and the Middle East alone, approximately US\$8 million is spent every year solely on removing infested palm trees, and in Italy, Spain, and France, RPW control and associated losses were projected to exceed US\$235 million by 2023 if effective containment programs were not implemented \cite{ashry2020}. By enabling early-stage detection before visible distress and significant structural damage occur, the proposed IoT-based acoustic network has the potential to dramatically reduce these costs by enabling timely curative chemical treatment rather than costly removal and replanting. At the farm level, the IoT-based networked acoustic architecture offers a scalable and cost-effective alternative to both invasive probe-based systems and high-cost optical or imaging technologies. Each node in the network consists of a low-cost acoustic microphone sensor paired with a wireless IoT communication module, enabling continuous real-time data transmission to a central processing unit without requiring human operators to be physically present at each tree \cite{karar2021}. The modular nature of the network means that farms of varying sizes can be instrumented incrementally, with additional sensor nodes deployed as needed, and that the system can be remotely managed and updated over-the-air. Saving even a fraction of the infested trees that would otherwise require removal yields immediate and measurable economic returns that offset the initial infrastructure investment. Furthermore, the automated nature of the network eliminates the recurring labor costs associated with periodic manual inspections, acoustic probe insertions, and visual surveys across large plantations \cite{saleem2022}.
 
\subsection{Societal Impact}
 
The societal implications of effective RPW detection extend well beyond the agricultural sector. Date palms hold deep cultural and economic significance across the Middle East, North Africa, and South Asia, where they have been cultivated for thousands of years and serve as a primary food source for millions of people. The fruit of the date palm provides essential nutrition, and its cultivation underpins the livelihoods of rural communities and smallholder farmers in these regions \cite{saleem2022}. RPW infestations threaten food security directly by destroying productive trees that require years or decades to mature and begin bearing fruit. Early detection preserves these long-maturing assets and protects the food supply chains that depend on them. For farming communities where date palm cultivation is the primary source of income, preventing infestation-related losses directly safeguards livelihoods, reduces rural poverty, and prevents farm abandonment \cite{saleem2022}. The IoT-based networked acoustic system offers particularly significant operational benefits for farm laborers and managers. Because each sensor node operates autonomously and transmits data wirelessly to a central hub, the system eliminates the need for repeated manual tree-by-tree inspections, a physically demanding and time-consuming task in large plantations, thereby reducing the labor burden on agricultural workers. Farm managers can monitor the infestation status of an entire plantation remotely through a connected interface, receiving automated alerts when any node detects acoustic signatures indicative of larval activity, without requiring constant on-site presence \cite{karar2021}. This shift from reactive manual inspection to proactive automated monitoring represents a fundamental improvement in the working conditions and operational efficiency of palm farming. From a public awareness and technology adoption perspective, the deployment of a connected IoT sensor network in agricultural settings represents a practical and accessible entry point into smart precision agriculture for farming communities in developing regions. Demonstrating the benefits of AI-driven, wirelessly networked pest detection in palm farming can serve as a model for adopting similar IoT-based monitoring frameworks in other crop protection contexts, fostering greater technological literacy and openness to innovation among farming communities.
 
\subsection{Global Impact}
 
RPW is a globally invasive pest species that has been detected in more than 60 countries across the Mediterranean basin, North Africa, the Middle East, South and Southeast Asia, and parts of the Caribbean and Central America \cite{saleem2022}. The rapid spread of RPW across continents reflects the interconnected nature of modern agricultural trade and the movement of plant material across borders. A scalable IoT-based acoustic detection network therefore has direct relevance to global biosecurity and pest containment efforts, as the wireless networked architecture can be replicated and deployed across geographically diverse farming environments with minimal adaptation. From an environmental perspective, the continuous 24/7 monitoring capability of an IoT acoustic network enables the targeted application of chemical treatments to individual infested trees identified by the system, rather than broad-scale prophylactic pesticide use across entire plantations. This precision approach made possible by the granular, per-tree visibility that the networked sensor architecture provides reduces the environmental burden of pesticide runoff, soil contamination, and harm to non-target organisms, aligning with global sustainability goals and the principles of integrated pest management (IPM). The reduction in unnecessary chemical use also decreases the risk of pesticide resistance developing in RPW populations, preserving the long-term efficacy of existing curative treatments \cite{saleem2022}. Finally, this project contributes to the global advancement of smart agriculture through the integration of IoT-networked bioacoustic sensing with automated deep learning classification. The system architecture developed --- encompassing wirelessly connected piezoelectric sensor nodes, on-node acoustic signal processing, feature extraction, and deep learning inference is transferable to other agricultural pest detection problems and biosecurity sensing applications in food security and environmental surveillance. Demonstrating a fully connected, remotely monitored IoT detection network operating in real farm conditions establishes a proof-of-concept that can accelerate the broader global adoption of intelligent, networked acoustic monitoring technologies in precision agriculture \cite{karar2021,alaufi2026}.

\chapter{Literature Review}
\label{ch:Literature}

% ===============================
% CHAPTER 2: LITERATURE REVIEW
% ===============================

\section{Acoustic Transduction Methodologies}
\label{sec:GRP-I}

Early detection of RPW infestation relies on vibroacoustic signals emitted by larvae as they burrow their way inside the palm tree trunk \cite{pinhas2008,karar2021}. During feeding and movement, larvae break the palm fibers, whose ``burst'' sounds are reported as impulsive sound trains \cite{pinhas2008,karar2021}. These signals serve as a major sign of infestation, since the internal tissue in a healthy palm tree tends to be acoustically quiet over the frequency range associated with the boring insects such as the RPW \cite{karar2021,sutanto2023}. Consequently, since visible external signs can appear only after severe internal damage has been done, acoustic emission is one of the most reliable indicators of active larval presence during earlier stages of infestation \cite{saleem2022,sutanto2023}.

\subsection{Vibroacoustic Sensing Instruments and Hardware Configuration}

Vibrational sensing instruments include piezoelectric crystals or contact microphones that are attached to the palm tree from which vibrations can be detected \cite{pinhas2008,karar2021}. Piezoelectric sensors convert mechanical stress of wood vibrations into measurable electrical signals by means of the piezoelectric effect \cite{karar2021,saranya2025}.

In contrast to conventional air microphones, which get affected and corrupted by environmental noise sources, contact probes are mechanically connected to the test trunk. In many implementations, a metal wave guide, for instance a stainless-steel rod or bolt, is disposed through a small pre-drilled hole in the trunk to improve acoustic coupling \cite{karar2021}. This arrangement enables the probe to operate in a similar manner to seismometer with internal feed vibrations being accentuated and airborne disturbance minimized \cite{karar2021,sutanto2023}.

\subsection{Signal Acquisition and Processing Techniques}

As larval signals are both amplitude and time structured, sophisticated signal processing is needed to separate them out from noise \cite{karar2021,saranya2025}. Raw signals are frequently sampled at 8 kHz to 44 kHz and then digital band-pass filtered, often in the range of 100 to 2500 Hz, to filter out low frequency disturbances like wind or machinery \cite{karar2022,pinhas2008,saranya2025}.

Feature extraction has been widely used to obtain compact spectral representations such as Mel-Frequency Cepstral Coefficients (MFCC), Bark Frequency Cepstral Coefficients (BFCC) and Linear Predictive Coding (LPC) \cite{saleem2022,saranya2025}. The Short-Time Fourier Transform (STFT) is often applied to convert time-domain signals into spectrogram images that visually represent energy distribution across frequencies \cite{karar2022,saranya2025}. These spectrograms are then used as inputs to machine learning classifiers, particularly Convolutional Neural Networks (CNNs), which learn to identify the distinct spectral-temporal patterns associated with larval feeding \cite{karar2022,saranya2025}. Architectures such as InceptionResNet-V2, MixConvNet, and Bidirectional Long Short-Term Memory (Bi-LSTM) networks have been explored for this purpose \cite{karar2021,karar2022,saranya2025}.

\subsection{Classification Performance and Experimental Accuracy}

Peer-reviewed studies show that sensitivity in controlled settings is high \cite{karar2021,karar2022,saranya2025}. Modified MixConvNet classifiers have reported accuracies of about 97.38\% on benchmark datasets \cite{karar2022}. Transfer-learning CNN pipelines using concatenated cepstral features have also achieved accuracies above 97\% on the public TreeVibes dataset \cite{boulila2023}. Combining LPC features with Bi-LSTM networks has achieved accuracies of 98.02\% \cite{saranya2025}. GMM/VQ-based statistical models reached detection rates close to 98.9\% in soundproofed laboratory experiments \cite{pinhas2008}. Those results illustrate the high capability of acoustic detection when combined with well-performed signal processing and classification schemes.

\subsection{Advantages and Operational Limitations of Acoustic Detection Systems}

Acoustic technology is especially valued in early detection; sensitive systems have reported detecting larval activity within roughly the first two weeks of infestation under controlled conditions \cite{ashry2020,saranya2025}. This early detection enables curative intervention before the structural and meristematic tissues of the palm are irreversibly damaged \cite{saleem2022,saranya2025}. Furthermore, the required hardware components, such as piezoelectric measuring heads and embedded IoT nodes, are generally much cheaper than high-end imaging technologies (for example, X-ray tomography and laser vibrometry), enabling affordable localized deployment of probe-based technologies.

However, there are several limitations hindering wide-spread deployment of probe-based acoustic systems. The need to drill in for the insertion of a waveguide introduces invasiveness and possibly provides entry points for other pests or pathogens \cite{karar2022}. Moreover, the detection radius of piezoelectric probes is usually about 1.5-2.0 m so that more than one sensor needs to be used for taller trees \cite{pinhas2008,karar2022}.

Environmental noise is still a challenge, and field studies have shown performance degradation under strong wind, heavy rain, and bird activity. From an operational point of view, the tools need to be installed and calibrated at each site, which may require substantial staff resources when setting up large commercial plantations \cite{karar2022,ashry2022cnn}.

\section{Distributed Acoustic Sensing (DAS) Methodologies} \label{sec:GRP-II}

Early detection of the Red Palm Weevil (RPW) relies on the vibroacoustic signals emitted by larvae as they feed and move through the internal tissue of a palm tree trunk \cite{ashry2022cnn,wang2021}. During feeding, larvae fracture palm fibers, generating impulsive burst sounds that propagate through the wood structure \cite{ashry2022cnn,wang2021}. These signals constitute a primary indicator of active infestation, since the internal tissue of a healthy palm tree is acoustically quiet across the frequency range associated with wood-boring insects such as the RPW \cite{pinhas2008,ashry2022cnn}. Because visible external symptoms of infestation typically appear only after severe internal structural damage has already occurred, acoustic emission detection represents one of the most reliable methods for identifying larval presence during the early stages of infestation \cite{pinhas2008,wang2021}.

\begin{figure}[H]
\centering
\includegraphics[width=0.75\textwidth]{figures/fig5_das_architecture.png}
\caption{DAS System Architecture \cite{wang2021}}
\label{fig:das_architecture}
\end{figure}

Unlike probe-based contact microphone systems, which are inserted invasively into individual tree trunks and offer only localized point sensing, Distributed Acoustic Sensing (DAS) based on phase-sensitive optical time-domain reflectometry ($\Phi$-OTDR) offers a fundamentally different sensing paradigm. In this approach, a single continuous optical fiber is non-invasively wound around the exterior of multiple tree trunks, allowing a single interrogation unit to monitor hundreds of trees simultaneously, and potentially on the order of 1,000 trees or more in larger deployments \cite{mao2021,ashry2020}. When larval feeding sounds perturb the refractive index of the fiber at the location of an infested tree, the resulting changes in the backscattered Rayleigh signal allow both the position of the disturbance and its frequency content to be identified in real time \cite{wang2021}.

Signal processing in DAS-based systems follows a pipeline broadly analogous to that used in contact probe systems. Raw time-domain Rayleigh traces are differenced against a reference to extract acoustic perturbations, and a band-pass filter spanning 100 to 800 Hz is applied to suppress low-frequency environmental noise such as tree swinging caused by wind and high-frequency electronic noise from the sensing hardware, while preserving the dominant spectral content of larval feeding sounds \cite{ashry2022cnn,wang2021}. The filtered signals are then organized as two-dimensional matrices of spatial points versus temporal or spectral samples and passed to convolutional neural network (CNN) classifiers trained through supervised learning to distinguish infested from healthy tree signatures \cite{ashry2022cnn,wang2021,mao2021}.

Classification accuracy across published DAS studies is consistently high under controlled and semi-controlled conditions. CNN models trained on temporal data achieved approximately 97.0\% accuracy, while models trained on spectral data achieved approximately 97.1\% accuracy in field experiments \cite{ashry2022cnn}. In laboratory environments designed to simulate farm conditions including wind and bird sounds as noise sources, CNN models achieved accuracies of up to 99.7\% on temporal data and 99.3\% on spectral data \cite{wang2021}. A merged classification approach, in which a signal is labelled as infested only when both the temporal and spectral CNN models independently agree, further reduced the false alarm rate by approximately 20\% relative to either model used alone \cite{ashry2022cnn}.

A particular advantage of the DAS methodology is its demonstrated capability for detecting larvae as young as 12 days old before any visible distress appears on the palm tree, enabling curative chemical treatment while the tree can still be saved \cite{ashry2020,ashry2022cnn}. The system operates continuously without human intervention, providing 24/7 monitoring coverage across large farm areas from a single sensing unit \cite{ashry2020}. However, the primary limitation identified in outdoor field deployments is the degradation of system performance at wind speeds exceeding approximately 9 mph, at which point wind-induced fiber vibrations produce acoustic signatures that overlap with larval feeding sounds and generate false positives \cite{ashry2022cnn}. Mitigating this wind sensitivity through thicker fiber jackets, improved signal processing, or adaptive noise cancellation remains an active area of investigation \cite{wang2021}.

\section{Non-Acoustic Detection Methodologies}
\label{sec:GRP-III}

This section examines alternative sensing methodologies beyond acoustic detection. Prior studies have investigated a range of approaches, including chemical, molecular, and pheromone-based sensing techniques. Evaluating these alternatives was essential to ensure that the most appropriate sensing methodology was selected for the proposed system.

\subsection{Volatile Organic Compound (VOC) Biosensing}

Chemical detection strategies for \textit{Rhynchophorus ferrugineus} rely on the identification and sensing of volatile organic compounds and other biochemical signatures associated with larval feeding and metabolic activity within date palm tissues. These approaches exploit the fact that infestation induces measurable chemical changes in the host tree and that the insect itself emits characteristic volatile compounds during development.

A whole-cell bacterial biosensor capable of detecting red palm weevil infestation in \textit{Phoenix dactylifera} was developed as a proof-of-concept system that translates specific chemical cues into measurable biological signals \cite{harpaz2022}. The biosensor is engineered to respond to infestation-associated metabolites, thereby enabling indirect detection of larval presence inside palm tissues. Since early-stage infestations are concealed within the trunk, visual inspection is typically ineffective; however, metabolic byproducts produced during larval feeding diffuse through plant tissues and can be detected externally. The biosensing platform utilizes genetically modified bacterial cells that produce a quantifiable output, such as fluorescence or colorimetric change, when exposed to target analytes linked to weevil activity \cite{harpaz2022}. This approach demonstrates that infestation-specific chemical markers can be transduced into measurable signals without destructive sampling of the tree, supporting its applicability for early detection scenarios.

\begin{figure}[H]
\centering
\includegraphics[width=0.99\textwidth]{figures/fig6_biosensor_schematic.png}
\caption{Schematic of The Experiment \cite{harpaz2022}}
\label{fig:biosensor_schematic}
\end{figure}

The chemical ecology of red palm weevil larvae further supports the feasibility of volatile-based detection. Detailed characterization of larval volatiles revealed stage-dependent emission profiles composed of specific chemical constituents that vary with developmental phase and pathogen stress \cite{ding2025}. The study identified distinct volatile organic compounds emitted by larvae and quantified their emission dynamics under varying biological conditions \cite{ding2025}. These findings indicate that larval metabolic processes generate a reproducible chemical fingerprint that can serve as a diagnostic indicator of infestation. Moreover, changes in volatile composition in response to physiological stress suggest that chemical emissions are sensitive to biological activity, potentially enabling differentiation between active and inactive infestations.

Integration of volatile profiling with biosensing technologies provides a pathway toward real-time chemical surveillance of date palms. Detection systems can be designed to monitor specific volatile compounds identified in larval emission studies, thereby enhancing selectivity and reducing false positives. The combination of targeted analyte identification and engineered biological sensing platforms establishes a chemically grounded framework for early detection of red palm weevil infestations \cite{harpaz2022,ding2025}. These methods enable non-invasive monitoring, are adaptable to field deployment, and leverage biochemical signatures intrinsic to insect presence and activity.

\subsection{Molecular Markers}

Molecular detection of \textit{Rhynchophorus ferrugineus} is based on the identification of species-specific genetic markers and ultrastructural traits that allow precise confirmation of infestation. Molecular techniques offer high specificity by targeting nucleic acid sequences unique to the red palm weevil, thereby enabling detection even when only trace biological material is present.

Molecular characterization and phylogenetic analysis of \textit{R. ferrugineus} populations in the Eastern Province of Saudi Arabia demonstrated the effectiveness of DNA-based identification using mitochondrial gene markers \cite{sabit2021}. Genetic sequencing revealed conserved regions suitable for species confirmation, while also identifying intra-species variation across populations \cite{sabit2021}. These results show that reliable species-level molecular identification is feasible and that conserved regions exist which could inform future polymerase chain reaction-based detection workflows using material recovered from infested trees.

Genetic variability and ultrastructural traits of adult red palm weevils from different geographical locations in Egypt were further examined using morphological and molecular techniques \cite{elzoghby2022}. The study demonstrated measurable genetic diversity among populations and documented diagnostic structural features that support taxonomic confirmation \cite{elzoghby2022}. Molecular markers derived from such analyses may help inform future rapid diagnostic assays, including species-specific primers for targeted amplification. The identification of stable genetic regions across populations can support detection robustness, while recognition of variability informs assay design to maintain sensitivity across different geographic strains.

Molecular detection methods offer several advantages for early-stage diagnosis. They enable confirmation at low infestation levels; they are highly specific; and they can differentiate red palm weevil from other sympatric insect species. By targeting conserved mitochondrial or nuclear gene regions identified in population studies, diagnostic protocols can be optimized for rapid screening of suspect trees \cite{sabit2021,elzoghby2022}. In principle, such assays could later be adapted to portable amplification systems for on-site testing. The integration of phylogenetic data with genetic variability assessments strengthens assay reliability and supports accurate detection across diverse ecological contexts.

\subsection{Pheromone Trap}

Pheromone-based detection exploits the aggregation behavior of \textit{Rhynchophorus ferrugineus}, which relies on chemical communication mediated by pheromones and synergistic attractants. By deploying optimized lure formulations and trap configurations, the presence of adult weevils in an area can be detected before widespread infestation occurs.

Optimization of pheromone trapping systems in Saudi Arabia demonstrated that trap performance is significantly influenced by lure composition, trap placement, and environmental factors \cite{shukla2018}. The study evaluated different trap contents and spatial arrangements, establishing conditions that maximize adult capture rates \cite{shukla2018}. Since adult weevils are responsible for dispersal and oviposition, their capture serves as an early indicator of infestation risk within a plantation. Effective trap optimization improves sensitivity, enabling detection of low-density populations.

Further assessment of attractant combinations in the United Arab Emirates examined the synergistic effects of pheromones with additional kairomonal components \cite{na2024}. The evaluation of different attractant mixtures demonstrated that specific combinations significantly enhance trap efficacy, increasing the probability of detecting active populations \cite{na2024}. Enhanced attractant formulations improve signal strength in surveillance programs by increasing adult aggregation around traps. This allows detection even when infestation levels are initially low.

\begin{figure}[H]
\centering
\includegraphics[width=0.2\textwidth]{figures/fig7_pheromone_traps.jpeg}
\caption{The Traps used in the Experiments \cite{na2024}}
\label{fig:pheromone_traps}
\end{figure}

Pheromone trap systems function as indirect detection mechanisms by capturing dispersing adults, thereby providing spatial information about infestation hotspots. Data collected from optimized traps can be used to map population distribution and monitor seasonal dynamics. The combined findings from trap optimization studies and attractant assessment experiments demonstrate that chemical communication pathways can be systematically leveraged for surveillance \cite{shukla2018,na2024}. Properly calibrated pheromone trapping networks enable early identification of infestation zones, facilitate targeted intervention, and support continuous monitoring of red palm weevil activity.

\section{Signal Processing Algorithms for Acoustic Data Processing of RPW}

The reliable detection of Red Palm Weevil (RPW) larvae through acoustic sensing depends not only on the quality of the hardware used to capture vibrational signals, but critically on the signal processing pipeline applied to raw acoustic data. Because larval feeding sounds are low-amplitude, intermittent, and embedded within a background of environmental noise, a sequence of carefully designed processing stages is required to extract meaningful features from the raw signals. The literature identifies a consistent pipeline consisting of preprocessing and filtering, followed by time-frequency feature extraction.

\subsection{Preprocessing and Bandpass Filtering}

Raw acoustic signals captured by piezoelectric contact microphones or similar transducers are typically sampled at rates ranging from 8 kHz to 44 kHz. Prior to any feature extraction, the signals undergo digital bandpass filtering to isolate the frequency of band of interest and suppress noise sources that fall outside the larval acoustic signature. The commonly applied bandpass range spans 100 to 2500 Hz, which encompasses the dominant chewing and movement sounds produced by RPW larvae while rejecting low-frequency interference such as wind-induced vibrations, ground tremors, and machinery noise below 100 Hz, as well as high-frequency electronic noise above the band of interest \cite{pinhas2008,rach2013}. This filtering stage is foundational to all subsequent processing steps, as it directly determines the signal-to-noise ratio of the features passed to any downstream analysis.

\subsection{Signal Squaring and Moving Average}

Following bandpass filtering, a common energy-based processing step involves squaring the time-domain signal to convert it into a strictly positive power envelope, which accentuates the amplitude differences between larval activity bursts and background noise. A moving average is then applied to the squared signal over a sliding time window typically 100 ms to smooth instantaneous noise spikes and reveal the underlying temporal pattern of larval feeding activity. This approach exploits the on/off feeding behavior of RPW larvae, which produces characteristic ups and downs in the moving average output when larvae are active, in contrast to the relatively flat and constant moving average observed from healthy trees. The complete signal processing pipeline, from raw signal acquisition through to the final infestation decision, is illustrated in Figure~\ref{fig:signal_pipeline} \cite{wang2021}.

\begin{figure}[H]
\centering
\includegraphics[width=0.99\textwidth]{figures/41598_2020_60171_Fig4_HTML.png}
\caption{Flowchart of the signal processing and detection algorithm \cite{ashry2020}}
\label{fig:signal_pipeline}
\end{figure}

\subsection{Signal-to-Noise Ratio Thresholding}

After computing the moving average, a Signal-to-Noise Ratio (SNR) is calculated for each measurement window, defined as the ratio between the high-level and low-level values of the moving average output. This definition is specifically chosen to mitigate false alarms caused by instantaneous noise spikes which affect only a brief portion of the window and therefore have limited impact on the moving average while reliably detecting the sustained amplitude modulation characteristic of larval feeding. If the computed SNR exceeds a predetermined decision threshold, the window is counted as an infestation alarm; otherwise, it is discarded. The decision threshold is selected empirically based on the SNR histogram distributions of infested and healthy tree measurements, and a threshold of 4 dB was found to yield 97 alarm counts for the infested tree versus only 9 for the healthy tree over a two-hour recording period \cite{wang2021}.

\subsection{Short-Time Fourier Transform and Spectrogram Representation}

A widely adopted alternative approach for converting time-domain acoustic signals into a form suitable for further analysis is the Short-Time Fourier Transform (STFT). The STFT divides the filtered signal into short, overlapping temporal windows and computes the Fourier transform of each window, yielding a two-dimensional time-frequency representation known as a spectrogram. The spectrogram encodes how the energy distribution across frequencies evolves over time, revealing the characteristic intermittent burst patterns associated with larval feeding \cite{pinhas2008,rach2013}. This representation is particularly valuable because it preserves both the spectral content and the temporal structure of the acoustic signal simultaneously, making the distinct impulsive nature of larval chewing sounds visually and mathematically distinguishable from more continuous environmental noise sources. Deep learning models applied to spectrogram inputs have demonstrated strong classification performance for RPW detection, with studies reporting accuracies exceeding 97\% under controlled conditions \cite{boulila2023}.

\subsection{Cepstral Feature Extraction}

An alternative and complementary approach to spectrogram-based representation is cepstral feature extraction. Mel-Frequency Cepstral Coefficients (MFCCs) compress the spectral content of an acoustic signal into a compact vector representation by mapping the power spectrum onto the mel frequency scale a perceptual scale that approximates the non-linear frequency response of the human auditory system and then applying a discrete cosine transform. Similarly, Bark Frequency Cepstral Coefficients (BFCCs) apply the same principle using the Bark psychoacoustic scale, and Constant Q Cepstral Coefficients (CQCCs) use a constant-Q frequency resolution that provides higher frequency resolution at lower frequencies. All three representations produce compact, noise-robust spectral descriptors that suppress redundant fine spectral detail while retaining the envelope shape most characteristic of the acoustic source \cite{pinhas2008,rach2013,boulila2023}. Studies have shown that concatenating MFCC, BFCC, and CQCC features before classification outperforms using any single feature extraction technique in isolation, yielding superior results across multiple acoustic datasets \cite{boulila2023}. Linear Predictive Coding (LPC) offers a further alternative, modeling the acoustic propagation channel through the wood structure as an autoregressive filter whose coefficients summarize the spectral envelope of the signal in a small set of parameters. LPC is computationally efficient and has been shown to effectively capture the resonance characteristics of larval feeding sounds as they propagate through palm trunk tissue \cite{saranya2025}.

\subsection{Comparative Assessment}

Each feature extraction approach involves trade-offs between computational complexity, compactness, and robustness to environmental noise. The SNR-based moving average approach is computationally lightweight and well-suited to real-time embedded deployment but is most effective when larval feeding produces clearly intermittent amplitude patterns and may be susceptible to false alarms under highly variable noise conditions \cite{wang2021}. STFT spectrograms preserve the most complete time-frequency information but produce large two-dimensional representations that carry higher computational cost in downstream processing \cite{boulila2023}. Cepstral features such as MFCCs, BFCCs, and CQCCs are significantly more compact and computationally lighter, and their combination has demonstrated particularly strong generalization across different acoustic environments \cite{boulila2023}. LPC coefficients are the most compact representation and least computationally demanding, making them attractive for embedded or IoT-constrained deployment scenarios \cite{saranya2025}. For real-world farm deployment, where wind, rain, bird activity, and mechanical noise are unavoidable, selecting the appropriate feature representation for the prevailing noise environment is a critical design decision that directly affects detection reliability \cite{pinhas2008,rach2013,saranya2025}.

\section{Deep Learning for Acoustic Data Classification of RPW Infestation}

\subsection{CNN Transfer Learning for Acoustic Signal Classification}

While BiLSTM captures temporal dependencies well, its performance relative to CNNs on smaller datasets is examined next. Many research papers have explored using transfer learning of pre-trained CNNs for acoustic signal classification. For example, the authors demonstrated that a fine-tuned transfer learning CNN classifier (InceptionResNet-V2) can be successfully applied to the TreeVibes dataset, achieving a classification accuracy of 94.53\% \cite{karar2021}. The TreeVibes dataset is an acoustic dataset of wood-feeding insects (such as RPW), where the data are represented as waveforms (time vs. relative amplitude) \cite{rigakis2021}.

In another study, the authors improved this approach by applying audio feature extractors to the raw audio signal before feeding the extracted features into transfer learning CNN models \cite{boulila2023}. The selected audio features were Mel-frequency cepstral coefficients (MFCC), Constant Q cepstral coefficients (CQCC), and Bark frequency cepstral coefficients (BFCC). Each of these feature extraction techniques generates a two-dimensional image representing the extracted features. These three feature images are concatenated and fed into one of the following transfer learning (TL) models: MobileNetV2, ResNet50V2, ResNet152V2, VGG16, VGG19, DenseNet121, DenseNet201, Xception, and InceptionV3.

The final feature vector output by a TL model is then fed into a classifier to determine whether the palm tree is infested. A research paper provided the following diagram to illustrate the full pipeline, from the audio signal to classification \cite{boulila2023}.

\begin{figure}[H]
\centering
\includegraphics[width=0.99\textwidth]{figures/fig9_cnn_pipeline.jpg}
\caption{Full Model Pipeline \cite{boulila2023}}
\label{fig:cnn_pipeline}
\end{figure}

When testing this architecture on the public TreeVibes dataset, an accuracy of 97.38\% was achieved \cite{boulila2023}. When testing on other environmental sound datasets, such as the ESC-10 dataset, the combination of MFCC, CQCC, and BFCC features yielded the best results. For ESC-10, this approach achieved ``the highest validation accuracy of 0.95 with a precision of 0.947, recall of 1.0, and an F1 score of 0.973'' \cite{boulila2023}.

These results show that combining MFCC, CQCC, and BFCC feature extractors with CNN-based transfer learning is a promising technique for detecting RPW infestations from audio signals.

\subsection{Bidirectional Long Short Term Memory Models for RPW Detection Through Acoustic Detection}

Another paper proposed a Bidirectional Long Short-Term Memory (BiLSTM) architecture for detecting RPW infestation from acoustic signals \cite{saranya2025}. The study introduced a complex, multistage data preprocessing pipeline, beginning with initial noise suppression to remove environmental noise such as wind. This is followed by Linear Predictive Coding (LPC), which extracts ``compact and discriminative'' spectral features from the audio signals that are specific to RPW activity.

Next, cosine similarity is computed between the extracted features and reference RPW and non-RPW vectors. The resulting cosine similarity coefficients are stored in a matrix, where each row corresponds to an audio file and each column contains similarity values. This matrix is then fed into a Bidirectional LSTM architecture, meaning that the temporal signal is processed by both a conventional forward LSTM and a backward (end-to-start) LSTM. The outputs of the forward and backward LSTMs are concatenated and passed to a dense layer for classification.

The researchers reported a classification accuracy of 98.02\%. They attributed the success of this approach in noisy environments to the ``combination of LPC feature extraction, cosine similarity-based feature matching, and a Bi-LSTM architecture.''

\subsection{Comparing RNNs and CNNs}

The researchers compared the use of CNNs and LSTM-based RNNs for environmental audio classification using the FSDKaggle2018 and ESC-50 datasets \cite{rezaul2024,fonseca2018}. Mel Frequency Cepstral Coefficients (MFCC) were used as a feature extractor, and the resulting feature vectors were fed into either a CNN or an LSTM model.

The ESC-50 CNN architecture used consisted of 10 layers, including convolutional layers with varying numbers of filters (ranging from 16 to 128), a filter size of 3, and ReLU activation, along with max-pooling, dropout, and dense layers for classification \cite{rezaul2024}. The LSTM architecture for ESC-50 consisted of 13 layers, including ``LSTM units, Batch Normalization, and Dropout layers.'' Time-distributed dense layers were also used, followed by a dense multi-class classifier layer.

The researchers reported that CNNs performed better than LSTMs on the smaller FSDKaggle2018 dataset, achieving 96.52\% and 87.88\% test accuracy, respectively \cite{rezaul2024,fonseca2018}. However, LSTMs performed better on the ESC-50 dataset, achieving 92.89\% accuracy compared to 88.47\% for CNNs \cite{rezaul2024}.

\subsection{CNN-LSTM Hybrid Architecture}

Another study proposed a hybrid CNN-LSTM model for environmental audio classification. This architecture uses 1D CNNs as feature extractors and feeds the extracted features into an LSTM block to capture time-based sequential relationships. A fully connected layer is then used for classification. Dropout layers are applied before the LSTM layer and before the fully connected layer. Max-pooling and ReLU activation are used after each convolutional layer. Figure~\ref{fig:cnn_lstm} illustrates the hybrid CNN-LSTM architecture presented in \cite{akter2025}.

\begin{figure}[H]
\centering
\includegraphics[width=0.85\textwidth]{figures/fig2_6_replacement.png}
\caption{CNN-LSTM Hybrid Model Architecture \cite{akter2025}}
\label{fig:cnn_lstm}
\end{figure}

This model achieved 99.33\% accuracy on the ESC-50 dataset and 98.38\% accuracy on the UrbanSound8K dataset. The researchers in \cite{akter2025} attributed this high accuracy to the use of MFCC to generate feature embeddings before feeding the data into the CNN-LSTM network, as well as incorporating ``batch normalization, transfer learning, and feature representation maps,'' which improved the ``convergence, generalization, and scalability'' of the model \cite{akter2025}. While these results are not specific to RPW detection, they show that combining CNN-LSTM hybrid architectures along with feature extraction techniques like MFCC could lead to incredible results for RPW detection. It is also promising to combine feature extraction techniques from \cite{boulila2023} (BFCC and CQCC) with the hybrid CNN-LSTM architecture proposed by \cite{akter2025}, as \cite{boulila2023} showed that combining MFCC, BFCC, CQCC feature extraction (through concatenation) produced better results than using any of these techniques individually.

\renewcommand{\arraystretch}{1.1}
\begin{landscape}
\scriptsize
\renewcommand{\arraystretch}{1.3}
\setlength{\tabcolsep}{6pt}

% Center the longtable on the page
\setlength\LTleft{\fill}
\setlength\LTright{\fill}

\begin{longtable}{@{}
    >{\raggedright\arraybackslash}p{2.8cm}
    >{\centering\arraybackslash}p{1.0cm}
    >{\raggedright\arraybackslash}p{1.8cm}
    >{\raggedright\arraybackslash}p{2.5cm}
    >{\raggedright\arraybackslash}p{9.5cm}
    >{\raggedright\arraybackslash}p{1.6cm}
    >{\raggedright\arraybackslash}p{2.0cm}@{}}
\caption{Summary of RPW Audio Classification Techniques adapted from the comparative review in \cite{boulila2023}}\label{tab:rpw_summary}\\
\toprule
\textbf{Ref.} & \textbf{Year} & \textbf{Algorithm} & \textbf{Data Type} & \textbf{Steps} & \textbf{Dataset} & \textbf{Accuracy} \\
\midrule
\endfirsthead
\caption[]{Summary of RPW Audio Classification Techniques (continued)}\\
\toprule
\textbf{Ref.} & \textbf{Year} & \textbf{Algorithm} & \textbf{Data Type} & \textbf{Steps} & \textbf{Dataset} & \textbf{Accuracy} \\
\midrule
\endhead
\bottomrule
\endfoot
\bottomrule
\endlastfoot

Pinhas et al. (2008) & 2008 & GMM, VQ & Spectral &
\textbullet\ Record the RPW using off-the-shelf devices \newline
\textbullet\ Manually label the data \newline
\textbullet\ Calculate the noise threshold \newline
\textbullet\ Count mismatched labels
& -- & 98.8\%, 98.9\% \\
\addlinespace

Parvathy et al. (2021) & 2021 & CAE & Mel Spectrograms &
\textbullet\ Capture audio signals \newline
\textbullet\ Extract features using the Mel spectrogram algorithm \newline
\textbullet\ Convert the features to images \newline
\textbullet\ Train the CAE on the images
& -- & 95.8\% \\
\addlinespace

Mohamed et al. (2021) & 2021 & AlexNet & MFCC &
\textbullet\ Extract features using the MFCC algorithm \newline
\textbullet\ Convert the features to images \newline
\textbullet\ Train AlexNet on the images
& -- & 99.2\% \\
\addlinespace

Wang et al. (2021) & 2021 & ANN, CNN & Temporal, Spectral &
\textbullet\ Install a fiber-optic DAS system on palm trees \newline
\textbullet\ Capture audio signals \newline
\textbullet\ Train ANN and CNN models with the DAS data \newline
\textbullet\ Deploy the DAS system in large farms
& -- & 99.9\%, 99.7\% \\
\addlinespace

Ashry et al. (2022) & 2022 & CNN & Temporal, Spectral &
\textbullet\ Record RPW sounds using DAS \newline
\textbullet\ Transform them into temporal and spectral data \newline
\textbullet\ Train CNN models on the data separately \newline
\textbullet\ Merge the CNN models
& -- & 97.0\%, 97.1\% \\
\addlinespace

Esmail Karar et al. (2022) & 2022 & CNN & Spectral &
\textbullet\ Extract spectral features \newline
\textbullet\ Convert the features to images \newline
\textbullet\ Train MixConvNet on the images \newline
\textbullet\ Compare performance with other models
& TreeVibes & 97.3\% \\

\end{longtable}
\end{landscape}

\section{Conclusions from Literature Survey}
\label{sec:Survey-Conclusion}

The literature presents that early detection of RPW infestation is technically feasible through acoustic, optical, chemical, molecular, and pheromone related techniques. However, these methods have tradeoffs among detection accuracy, invasiveness, scalability, and the cost of deployment.

Acoustic probe-based methods have demonstrated good potential for early detection, with machine learning classifiers reaching over 97\% accuracy in laboratory settings. The systems can capture larval feeding vibrations in the 100--2500 Hz range and can detect infestation within roughly the first two weeks under controlled conditions. However, their invasive installation requirements, limited sensing radius (1.5--2.0 m), and sensitivity to environmental noise significantly restrict large-scale deployment in commercial plantations.

DAS offers a scalable and non-invasive alternative, enabling simultaneous monitoring of many trees using fiber-optic cables. When integrated with CNN-based classifiers, the proposed DAS systems showed high detection performance and low false alarm rates. Yet fiber costs, interrogation head cost and installation are still too high for widespread use, especially in the price sensitive agricultural sector.

Tasks using chemical biosensing and molecular markers are highly specific to infestation and can identify it through metabolic or genetic changes. Although these are potential approaches for accurately describing biological activities, many of them involve processing in the laboratory environment or use reagents that need to be tailored to meet specificity requirements and be optimized across a bioengineering platform, not feasible when considering continuous real-time surveillance over large farms.

Pheromone trapping systems serve as useful surveillance tools for adult weevil populations but cannot detect internal larval activity. Since structural damage occurs during the larval stage, reliance on adult capture alone does not guarantee early-stage intervention.

A common limitation across all approaches is that, although high detection accuracy is attainable under controlled conditions, environmental noise (i.e., field deployment), scalability (i.e., the cost of equipment and labor), invasiveness, and operational complexity continue to create major challenges. Even though deep learning architectures (e.g., MixConvNet, InceptionResNet-V2 and Bi-LSTM) have recently shown competitive classification accuracy, their performance in the diverse outdoor environment remains dataset-specific and computationally intensive.

It is clear that there is an unmet need for a low-cost, non-invasive, wearable IoT-based acoustic detection system that strikes a balance between affordability and intelligent signal classification. Unlike fiber-optic DAS systems, which are currently deployed along the trunk, the proposed low-cost piezoelectric trunk-mounted sensors and cloud-based deep learning pipeline offer low infrastructure costs and automated classification capability.

The most significant contribution of this work, though, is its focus on practical scalability: developing a modular wireless sensor network architecture that can be deployed for large scale plantations without the excessive cost and installation effort traditionally associated with optical fiber-based systems. Combining low-cost hardware with machine learning-based acoustic discrimination, the project intends to bring laboratory-quality accuracy in acoustics closer to on-farm applicability.

\chapter{System Specifications}
\label{ch:Specifications}
 
% ===============================
% CHAPTER 3: SYSTEM SPECIFICATION
% ===============================
 
\section{Functional Requirements}
\label{sec:FR}
 
FR1: The system must detect internal palm tree vibrations using a piezoelectric sensor interfaced with an ESP32 sensing node and convert them into digital signals for processing.
 
FR2: The system must process the captured acoustic signals on the sensing node by per- forming signal conditioning, SNR-based filtering, and MFCC feature extraction.
 
FR3: The system must transmit only significant acoustic events, determined by an SNR threshold, from sensing nodes to the cluster head using LoRa communication along with node ID and timestamp.
 
FR4: The system must aggregate incoming data at the cluster head and reliably forward it to the cloud backend using MQTT over available network interfaces (Wi-Fi or cellular).
 
FR5: The system must analyze incoming acoustic feature data in the cloud backend using a trained machine learning model to classify and detect potential red palm weevil infestations.
 
FR6: The system must store acoustic features, classification results, node health, and cluster health data in a centralized time-series database.
 
FR7: The system must generate and log alerts when infestation events are detected and notify users through automated mechanisms such as email.
 
FR8: The system must provide a web-based dashboard that visualizes sensor node status, geographic distribution, detection events, and historical trends.
 
FR9: The system must allow authenticated users to securely access system data, including real-time monitoring and historical records, through the dashboard.

\section{Non-Functional Requirements}
\label{sec:NFR}
 
NFR1: The system should achieve reliable detection of red palm weevil activity by capturing and processing acoustic signals within the 100 Hz to 2500 Hz frequency range.
 
NFR2: The system should optimize bandwidth usage and energy consumption by transmitting only processed feature data instead of continuous raw audio streams.
 
NFR3: The system should support large-scale deployment, enabling thousands of sensing nodes to operate under a single cluster head without significant performance degradation.
 
NFR4: The system should ensure low-latency and reliable communication between sensing nodes, cluster heads, and the cloud backend, even under intermittent connectivity conditions.
 
NFR5: The system should ensure secure data transmission and access control using encryption protocols and authentication mechanisms.
 
NFR6: The system should efficiently store and query large volumes of time-series data to support real-time monitoring and long-term analysis.
 
NFR7: The system should provide a responsive, intuitive, and user-friendly web interface accessible across standard desktop and tablet browsers.
 
NFR8: The system should be robust for outdoor deployment, ensuring hardware durability and stable operation under environmental conditions such as heat, humidity, and dust.
 
\section{External Interfaces and Database Requirements}
 
\subsection*{External Interfaces}
 
The RPW acoustic detection system interacts with external components spanning three distinct layers: the sensing layer, the cluster aggregation layer, and the cloud backend.
 
\subsubsection*{Sensor-to-Microcontroller Interface}
 
At each sensing node, the primary hardware interface is between the 27mm piezoelectric disc sensor and the ESP32 microcontroller. The piezoelectric element is in physical contact with the surface of the palm trunk and generates a low-amplitude analog voltage signal (mV range) that varies as a function of internal vibrations generated by larval feeding activity. The output signal is then fed into an MCP6002 op-amp amplification circuit, which brings the output value to a level suitable for ESP32's analog to digital converter (ADC) input range. The conditioned signal is then digitized by an ADC at a rate sufficient to capture the full RPW larval frequency band, known to range from 100 Hz to 2500 Hz \cite{pinhas2008,rach2013}, with a desired sampling rate of at least 8 kHz (i.e., well above the Nyquist limit).
 
\subsubsection*{Sensing Node-to-Cluster Head Interface}
 
Communication between the sensing nodes and the cluster head is established using Arduino's HopeRF RFM95W module at SF7 through LoRa. Transmissions are coordinated with a time division multiple access (TDMA) slot scheme with 900 ms time slots allocated over a 24-hour period across two biologically informed transmission windows: nighttime (10 PM--6 AM; hourly wake-ups) and daytime (6 AM--10 PM; four-hour intervals). All transmissions consist of a small payload of 214 bytes that includes the node identifier, a UTC timestamp, and a pre-computed MFCC feature vector. No raw audio is ever sent; feature extraction occurs on the sensing node, and only features are transmitted. A transmission event is generated only when the Signal-to-Noise Ratio (SNR) exceeds 4 dB; otherwise, the transmission is suppressed and bandwidth is conserved \cite{wang2021}. With this setup, one cluster head operating at SF7 can support 8,000 sensing nodes, meaning that LoRa TDMA capacity, rather than radio range, becomes the limiting factor for network size.
 
\subsubsection*{Cluster Head-to-Cloud Interface}
 
A Raspberry Pi 4 + RFM95W LoRa module is used for downlink reception, building on each cluster head. The cluster head runs continuously on mains power, receiving over both transmission windows and writing all arrived MFCC feature packets to a local SD card. Wi-Fi (IEEE 802.11 b/g/n) is used to forward packets received from the meter node level to the centralized cloud backend where possible, and for remote display heads where only a 4G cellular uplink is available. The communication between the nodes takes place through data sent via the MQTT protocol over TCP/IP, and each message contains a JSON-encoded payload with the node identifier, UTC timestamp, and MFCC feature vector. All incoming messages are sent to and routed by the cloud broker to respective processing and storage pipelines. Classification is not performed on the cluster head, which acts only as an always-on aggregation and forwarding layer. In the case of large plantations, many cluster heads are used in non-overlapping geographic zones with each designated a unique block of Node IDs. No inter-cluster-head coordination is required. The system scales linearly, adding extra network capacity with each new cluster head in the second layer (8,000 nodes block per cluster head) and the estimated LoRa radio range is roughly equal to 1 km for every deployed cluster head in a real plantation scenario.
 
\subsubsection*{Cloud-to-Dashboard Interface}
 
The Cloud-to-Dashboard Interface is a web-based monitoring dashboard, and the cloud backend exposes a RESTful HTTP API to serve this data. Plantation managers access the dashboard through a standard web browser, which communicates with the backend over HTTPS at port 443. The API endpoints can be used to obtain sensor status, historical (MFCC) feature records, classification results and alert logs. To make sure plantation data is accessed only by authorized parties, authentication is processed using JSON Web Tokens (JWT). Automated email alerts are dispatched from the cloud server to plantation manager accounts after detected infestation events.
 
\subsubsection*{User Interface}
 
The web dashboard is the main interface between the system and its end users. The single page application (SPA) implements dynamic visualization of sensor node and cluster head locations on a real-time map with color-coded insect infestation status markers along with time-series graphs for acoustic energy levels by node, timestamps and confidence scores against detection events, as well as system health indicators such as connectivity between all nodes in the network and battery levels. The web dashboard is accessible through any desktop and tablet browser with no software to install.
 
\subsection*{Database Requirements}
 
We use a PostgreSQL database with the TimescaleDB extension for our primary data store. PostgreSQL offers a robust, fully-supported relational database engine that allows for querying via standard SQL. TimescaleDB is an extension of PostgreSQL that provides automatic time-stratified partitioning and tuned range-query performance, which suits the continuous time-stamped MFCC feature streams output by the sensing layer. This combination also manages relational relationships of the system, metadata for cluster heads and nodes, and user accounts in a single database instance; no separate data store is required.
 
\subsubsection*{Logical Database Organization}
 
The database is organized into the following tables, each recording timestamped entries keyed by node identifier and farm zone:
 
\begin{table}[H]
\centering
\caption{Database Table Index}
\label{tab:db_schema}
\renewcommand{\arraystretch}{1.3}
\small
\setlength{\tabcolsep}{4pt}
\begin{tabular}{>{\raggedright\arraybackslash}p{0.24\textwidth} >{\raggedright\arraybackslash}p{0.15\textwidth} >{\raggedright\arraybackslash}p{0.15\textwidth} >{\raggedright\arraybackslash}p{0.40\textwidth}}
\toprule
\textbf{Table} & \textbf{Index} & \textbf{Columns} & \textbf{Data Columns} \\
\midrule
\multicolumn{4}{l}{\textbf{Sensor Layer}} \\
nodes & node\_id & cluster\_id (varchar), & farm\_zone (varchar), install\_date (date), \\
      &          &                        & location\_lat (float), location\_lon (float) \\
node\_health & timestamp, node\_id & & battery\_mv (int), rssi\_dbm (int), uptime\_s (int) \\
acoustic\_features & timestamp, node\_id & & mfcc (float[]), snr\_db (float) \\
\midrule
\multicolumn{4}{l}{\textbf{Cluster Layer}} \\
cluster\_heads & cluster\_id & & farm\_zone (varchar), install\_date (date), location\_lat (float), location\_lon (float), uplink\_type (varchar) \\
cluster\_health & timestamp, cluster\_id & & uptime\_s (int), rssi\_dbm (int), sd\_free\_mb (int), uplink\_type (varchar) \\
\midrule
\multicolumn{4}{l}{\textbf{Cloud Layer}} \\
classification\_results & timestamp, node\_id & & label (varchar), confidence (float), model\_version (varchar) \\
alert\_log & timestamp, node\_id & & alert\_type (varchar), confidence (float), acknowledged (bool) \\
\bottomrule
\end{tabular}
\end{table}
 
The acoustic features, classification results, cluster health, and node health tables are declared as TimescaleDB hypertables, enabling automatic time-based partitioning and accelerated time-range queries. Raw acoustic data is no longer retained at the database level, as feature extraction is performed on-node prior to transmission. Extracted MFCC features and classification results are retained for one year. Alert logs and node and cluster head metadata are retained indefinitely to support long-term infestation trend analysis and system auditing.

\chapter{Technical Approach}
\label{ch:TechnicalApproach}

% ===============================
% CHAPTER 4: Technical Approach
% ===============================

\section{Proposed Solution}
\label{sec:ProposedSolution}

In order to overcome the drawbacks of existing RPW detection and monitoring methods, this study proposes a low-cost 1D data acquisition method using piezoelectric acoustic sensors. This data is then transmitted wirelessly to an aggregate layer, followed by another wireless transmission to a cloud server, where the data is analyzed and classified using deep learning models. The system uses arrays of trunk-mounted piezoelectric sensors interfaced with ESP32 microcontrollers all communicating over a LoRa-based wireless sensor network to a centralized cloud backend. This architecture is constructed to provide high levels of detection precision at a low cost compared to alternatives, allowing scalable activity across large commercial plantations.

The system consists of a three-layered architecture, with each layer handling a specific step during the detection pipeline; these are Sensing Layer, Cluster Layer, and Cloud Layer. In Figure 4.1, the high-level block diagram for the proposed solution is shown.

\begin{figure}[h]
\centering
\includegraphics[width=1\textwidth]{figures/system_fig_diagram.pdf}
\caption{High-level block diagram of the proposed RPW acoustic detection system}
\label{fig:HighLevelBlockDiagram}
\end{figure}

\subsection{Sensing Layer}

The Sensing Layer is the field-deployed layer of the system. Each sensing node is composed of a piezoelectric disc sensor mechanically coupled to the surface of a palm trunk, an op-amp for analog signal conditioning, an ESP32 microcontroller for digitization and on-node processing, and a LoRa transceiver for wireless transmission. The piezoelectric element translates the mechanical vibrations produced by feeding activity of larvae resident inside the trunk to low-amplitude analog voltage signals in the millivolt range. All signal processing is performed on the ESP32 before any data leaves the sensing node. No raw audio is transmitted; only the extracted MFCC vector, along with the node identifier and a UTC timestamp, is packaged into a LoRa payload and sent to the cluster layer. The sensing nodes operate on a duty-cycled schedule aligned with the biological activity patterns of RPW larvae. This schedule maximizes detection probability while reducing power consumption, allowing nodes to run on batteries for extended periods with optional solar charging.

\subsection{Cluster Layer}

The Cluster Layer is the always-on aggregation and forwarding tier between the field-deployed sensing nodes and the cloud backend. Each cluster head consists of a Raspberry Pi 4 Model B equipped with an RFM95W LoRa receiver module interfaced via SPI. The cluster head runs continuously on mains electricity and remains in continuous receive mode, listening for incoming LoRa packets from sensing nodes.

Upon receiving a LoRa packet, the cluster head validates the packet header and CRC check-sum. Valid packets are parsed to extract the a 2-byte node ID, 4-byte timestamp, 1-byte SNR, the compressed MFCC representation, and four burst-statistic features. The parsed data is written to local storage (SQLite on an SD card) as a buffer against connectivity interruptions. If an uplink connection is available, whether Wi-Fi (IEEE 802.11 b/g/n) or 4G cellular, the cluster head publishes the data to the cloud backend via MQTT over TCP/IP as a JSON-encoded payload. If the uplink is temporarily unavailable, the data is buffered locally and forwarded upon reconnection. The cluster head does not perform classification or inference; it only aggregates and forwards data.

For large plantations, multiple cluster heads are deployed in non-overlapping geographic zones, each assigned a unique block of node identifiers. No inter-cluster-head coordination is required; the system scales linearly by adding cluster heads.


\subsection{Cloud Layer}

The Cloud Layer provides centralized data storage, deep learning inference, alert generation, and a monitoring interface. The cloud backend receives MQTT messages from all cluster heads, routes them through a message broker, and stores the incoming MFCC feature data alongside node and cluster health metadata in a PostgreSQL database extended with TimescaleDB for time-series storage and querying. Incoming acoustic feature data is classified by a trained deep learning model, based on CNN, CNN-LSTM, or Mamba (SSM) architectures, to determine whether the acoustic signature corresponds to active larval feeding or environmental noise. When the classification confidence exceeds a predefined threshold, the system generates an alert, logs it in the alert database, and sends automated email notifications to registered plantation managers. The web-based monitoring dashboard is the primary user interface, providing real-time visualization of sensor nodes, detection events, and system health. The dashboard communicates with the cloud backend through a RESTful HTTP API secured with HTTPS (TLS 1.2 or later) and JWT authentication, so plantation data is accessible only to authorized users. It runs in any standard desktop or tablet browser with no installation required.

\section{Hardware Design}

The hardware architecture of the system is organized into two layers: the sensing node (deployed per tree) and the cluster head (deployed per farm zone), with data ultimately flowing to a cloud backend for classification and monitoring. Below is a detailed description of every component in the system.

\subsection{Piezoelectric Disc Sensor}

Each sensing node begins with a 27 mm piezoelectric disc (Murata 7BB-27-4L0 or equivalent) that is placed directly on the surface of the palm tree trunk. The piezoelectric element converts mechanical vibrations produced by larval feeding activity into low-amplitude analog voltage signals in the millivolt range. This non-invasive surface-contact approach was chosen over air-coupled microphones because it is far less susceptible to ambient environmental noise such as wind, bird calls, and machinery, and it focuses specifically on vibrations propagating through the wood structure \cite{pinhas2008,karar2022}.

\subsection{MCP6002 Op-Amp Signal Conditioning Circuit}

The raw output of the piezoelectric sensor is too weak to be sampled reliably by the microcontroller's ADC. An MCP6002 operational amplifier circuit amplifies and conditions the signal, bringing it into the 0 to 3.3 V input range expected by the ESP32. Protective components including 1N4148 clamping diodes, decoupling capacitors, and a voltage divider network ensure that voltage spikes do not damage the ADC input. This stage also contributes to the initial noise rejection by shaping the signal bandwidth before digital filtering takes place.

\subsection{ESP32 Microcontroller}

The ESP32 serves as the brain of the sensing node. It performs ADC sampling and executes the on-node signal processing pipeline described in Section 4.3.3. No raw audio ever leaves the node; only the pre-computed features are transmitted, which dramatically reduces bandwidth and power consumption. The ESP32 was selected for its low cost, integrated Wi-Fi (useful for bench testing), adequate ADC resolution, and broad community support through the Arduino ecosystem.

\subsection{LoRa Transceivers}

Wireless communication between the sensing node and the cluster head is handled by the RFM95W, which interfaces with the ESP32 over SPI. Operating at 915 MHz in the license-free ISM band at spreading factor 7, it provides a practical radio range of approximately 1 km per cluster head in a real plantation scenario. Transmission scheduling and payload composition are described in Section 4.3.3. At SF7, a single channel can theoretically support approximately 8,000 measurement transmissions per hour under the defined payload and airtime constraints.

\subsection{LoRa Antenna}

Both the sensing node and the cluster head use dedicated 915 MHz RPSMA antennas to maximize transmission range and link reliability. An external antenna, as opposed to an onboard PCB trace antenna, was chosen because the nodes will be mounted inside weatherproof enclosures on tree trunks, where RF shielding from the casing and the trunk itself could degrade signal quality.

\subsection{Raspberry Pi 4 Model B}

The cluster head runs on a Raspberry Pi 4 with 4 GB RAM, powered continuously by a standard power outlet. It receives incoming LoRa packets from all sensing nodes in its zone via a second RFM95W LoRa module connected over SPI, logs each received MFCC feature packet to a local 32 GB SD card for redundancy, and forwards the data to the cloud backend over Wi-Fi (IEEE 802.11 b/g/n) or, where Wi-Fi is unavailable, data will be stored in a buffer until uplink is available. The cluster head does not perform any classification; it functions purely as an always-on aggregation and forwarding layer. The Raspberry Pi OS (Debian-based, 64-bit) runs a Python 3.10 daemon that handles packet reception, local logging, and MQTT publishing.

\subsection{Casing \& Optional Solar Panel}

Each node is enclosed in an IP65-rated weatherproof casing and optionally powered by a solar panel with battery regulation circuit, enabling autonomous outdoor operation without mains connectivity.


\section{Software Design}

This section presents the software design of the RPW acoustic detection system. It covers the overall system architecture, the data architecture underlying persistent storage and data flow, and the detailed design of both static and dynamic system behaviors. Each diagram included in this section is discussed in the accompanying text.



\subsection{System Architecture}

A detailed description of the overall physical architecture of the application is presented in this section.

\subsubsection{Node Architecture}

The IoT-based RPW acoustic detection system is organized into a three-tier physical architecture comprising (1) the sensing node layer, (2) the cluster head aggregation layer, and (3) the cloud backend and dashboard layer. Each tier has clearly defined responsibilities and communicates with adjacent tiers through well-specified interfaces, enabling modular development, independent scaling, and fault isolation across the deployed network.

\paragraph{Sensing Node Layer} 

At the lowest tier, each sensing node uses the hardware described in Section 4.2: a piezoelectric disc sensor coupled to the palm trunk, an MCP6002 op-amp conditioning circuit, an ESP32-WROOM-32 microcontroller, and an RFM95W LoRa transceiver. The piezoelectric element converts internal mechanical vibrations generated by RPW larval feeding into a low-amplitude analog voltage signal, which is conditioned by the op-amp circuit before being digitized by the ESP32's onboard ADC at a sampling rate of at least 8 kHz, which satisfies the Nyquist criterion for the larval acoustic band of 100--2500 Hz. Once sampled, the ESP32 firmware applies a digital bandpass filter (100--2500 Hz), squares the filtered signal, computes a moving average to produce a power envelope, and calculates the Signal-to-Noise Ratio (SNR) of each measurement window. A transmission event is generated only when the SNR exceeds a 4 dB threshold, ensuring that bandwidth and energy are consumed only for acoustically significant events. When triggered, the node extracts a Mel-Frequency Cepstral Coefficient (MFCC) feature vector on-device and transmits a compact ~214-byte payload consisting of a 15×13 int8-quantized MFCC representation with a global scale factor, 4 burst-statistic scalars, and 7 bytes of metadata (node ID, timestamp, SNR) to the cluster head via the LoRa transceiver operating at SF7. No raw audio is ever transmitted across the network. Transmissions are coordinated using a Time Division Multiple Access (TDMA) scheme with 900 ms time slots distributed across two biologically informed windows: night-time (22:00--06:00, hourly) and daytime (06:00--22:00, four-hour intervals). This schedule aligns monitoring frequency with periods of heightened larval acoustic activity.

\paragraph{Cluster Head Aggregation Layer}

The intermediate tier consists of one or more cluster heads, each built around a Raspberry Pi 4 Model B paired with an RFM95W LoRa receive module. The cluster head operates continuously on mains power and listens across both defined transmission windows, receiving MFCC feature packets from all sensing nodes within its LoRa coverage zone (approximately 1 km radius in open plantation conditions). All received packets are written to a local 32 GB microSD card for store-and-forward resilience, then forwarded to the cloud backend using MQTT v3 over TCP/IP. Where Wi-Fi coverage (IEEE 802.11 b/g/n) is available, that uplink is used; in remote zones lacking Wi-Fi, the cluster head falls back to a 4G cellular uplink conforming to 3GPP LTE standards. Each MQTT message carries a JSON-encoded payload containing the node ID, UTC timestamp, SNR, a 15×13 int8 MFCC representation with scale factor, and four burst-statistic features (burst count, mean inter-burst interval, peak SNR, spectral centroid). The cluster head performs no classification; it serves exclusively as an always-on aggregation and forwarding tier. Because a single cluster head operating at SF7 can support up to 8,000 sensing nodes within a TDMA frame, the limiting factor for network size is LoRa TDMA capacity rather than radio range, allowing linear scaling by simply deploying additional cluster heads in non-overlapping geographic blocks.

\paragraph{Cloud Backend and Dashboard Layer}

At the top tier, a Linux-based cloud server (Ubuntu 22.04 LTS) receives all MQTT messages via a cloud MQTT broker and routes them through an automated preprocessing pipeline implemented in Python 3.10. Incoming MFCC feature vectors are validated, normalized, and passed to a trained deep learning inference service, which classifies each acoustic event as either RPW larval activity or ambient noise. Classification results, confidence scores, raw MFCC features, and node health telemetry are stored in a PostgreSQL database extended with TimescaleDB, which provides automatic time-based partitioning and accelerated range-query performance for the continuous time-stamped data streams. When a classification result exceeds a predefined confidence threshold, the alert engine automatically dispatches an email notification to registered plantation managers, identifying the affected node ID, farm zone, and detection confidence. A RESTful HTTP API (HTTPS, port 443, JWT-authenticated) exposes sensor status, historical feature records, classification results, and alert logs to the web-based monitoring dashboard. The single-page dashboard application provides real-time visualization of sensor node and cluster head locations on a map with color-coded infestation status markers, time-series graphs of acoustic energy levels per node, detection event timestamps with confidence scores, and system health indicators including battery levels and connectivity status for all nodes in the network.

\subsubsection{Node Topology}

The node architecture adopted for this system is a two-level hierarchical wireless sensor network (WSN). Sensing nodes communicate exclusively with their designated cluster head using LoRa (long-range, low-power) radio in the license-free ISM band at 915 MHz. This star-of-stars topology was selected over a flat single-hop topology because it concentrates the energy cost of long-distance transmission at the always-on, mains-powered cluster head rather than at battery-constrained sensing nodes, and over a mesh topology because it eliminates the protocol overhead and latency introduced by multi-hop routing while still achieving the required plantation-wide coverage. The cluster head then forwards aggregated data to the cloud over a high-bandwidth backhaul (Wi-Fi or 4G), entirely decoupling the constrained sensing layer from the cloud infrastructure. This architecture is illustrated conceptually in Figure 4.2 below.


\begin{figure}[h]
\centering
\includegraphics[width=0.95\textwidth]{figures/3Node_Arch.png}
\caption{Three-Tier Node Architecture of the IoT-Based RPW Acoustic Detection System}
\label{fig:ThreeTierNodeArchitecture}
\end{figure}


\subsubsection{Component Architecture}

Each tier of the system is composed of several internal software components that handle distinct stages of the detection pipeline. This section describes the responsibilities of each component and how they are organized within their respective nodes.

\paragraph{Sensing Node}

The sleep manager governs the ESP32's deep sleep and wake cycle, triggering wake-ups at the intervals defined by the active transmission window and returning the MCU to low-power mode after each recording and processing cycle completes. The ADC component captures a fixed-duration audio sample from the conditioned piezoelectric input each time the node wakes, buffering it in memory for downstream processing. The signal processing module runs the full on-device analysis pipeline, from bandpass filtering through power envelope computation and SNR evaluation to MFCC extraction, and makes the transmit-or-discard decision based on the SNR outcome. The LoRa transmitter handles payload assembly, TDMA slot timing, and packet dispatch through the RFM95W module. The power manager monitors battery voltage, manages solar charging state, and embeds a low-battery flag in outgoing packets when the node is approaching depletion so that operators can schedule maintenance before data loss occurs.

\paragraph{Cluster Head}

The LoRa listener is a continuously running daemon that receives incoming packets from sensing nodes, validates packet headers and CRC checksums, and discards malformed transmissions before passing valid data to the next stage. The storage component writes every validated packet to the local SD card as a timestamped entry, acting as a persistent buffer that retains data independently of uplink availability. The MQTT publisher reads from this buffer and forwards packets to the cloud broker whenever connectivity is present, automatically retransmitting any entries that were held during a connectivity interruption. The health manager periodically collects system-level diagnostics, including SD card free space, daemon uptime, and average RSSI across connected sensing nodes, and publishes them alongside the acoustic data so that infrastructure health is visible through the dashboard.

\paragraph{Cloud Backend}

The MQTT broker receives all incoming messages from cluster heads and routes them into the processing pipeline. The preprocessing service validates and normalizes incoming MFCC feature vectors before passing them to the inference engine. The inference engine loads the trained deep learning model and classifies each acoustic event, producing a label and a confidence score that are written to the database alongside the original features. The alert engine monitors classification outputs and, when the confidence score exceeds the configured threshold, generates a log entry and dispatches an email notification identifying the affected node and farm zone. The REST API serves stored data, including sensor status, classification history, and alert logs, to authenticated clients over HTTPS. The dashboard frontend consumes the API and renders the monitoring interface, providing the real-time map view, time-series graphs, detection event logs, and system health indicators that plantation managers interact with.

\subsubsection{Class Architecture and Static Representation}

The class architecture consists of 4 classes: ClusterHead, SensingNode, SensorMeasurement, and Tree. The ClusterHead and SensingNode classes represent the hardware components of Cluster Head and Sensing Node. Each ClusterHead can have many Sensing nodes connected to it, and each SensingNode can be connected to one ClusterHead. This relationship is represented by a bidirectional aggregation association relationship.


\begin{figure}[h]
\centering
\includegraphics[width=1\textwidth,height=1\textheight,keepaspectratio]{uploads/Class Daigram.png}
\caption{Class architecture diagram of the RPW detection system}
\label{fig:ClassArchitectureDiagram}
\end{figure}

The SensorMeasurement class represents one “Chunk” of acoustic measurements, with a start time stamp, end time stamp, and an acoustic payload. Each SensingNode can have many SensorMeasurement stored in a dynamic array (vector). Each SensorMeasurement can only belong to one SensingNode. This is represented by a 1 to many composition relationship, because without a SensingNode no SensorMeasurement can be made. Each SensingNode is also associated with one tree, and each tree is associated with one sensing node, this relationship is represented by a one-to-one composition association, because without a SensingNode we cannot monitor a tree, and therefore for a tree to exist in our system it must be associated with exactly one sensing node.

\subsection{Data Architecture}

The data architecture is built on PostgreSQL extended with TimescaleDB. PostgreSQL provides a relational database engine supporting complex queries via standard SQL, while TimescaleDB adds automatic time-based partitioning (hypertables) and range-query performance suited to the continuous timestamped MFCC feature streams produced by the sensing layer. Together, they manage all relational data, cluster head and node metadata, time-series acoustic features, classification results, and user accounts, within a single database instance. The database schema consists of seven tables distributed across the three system layers. Figure 4.4 presents the ER diagram of the complete schema.

\subsubsection{Sensing Layer Tables}

The nodes table stores the registration record for each deployed sensing node, identified by a node id primary key (varchar). Each row contains the associated cluster id foreign key, the farm zone designation, geographic coordinates (location lat, location lon), and install date.

\begin{figure}[h]
\centering
\includegraphics[width=0.95\textwidth]{figures/ER_Diagram.pdf}
\caption{Entity-Relationship diagram of the RPW detection system database}
\label{fig:EntityRelationshipDiagram}
\end{figure}

The relationship to cluster heads is one-to-many: each cluster head may have many sensing nodes, but each node belongs to exactly one cluster head. The node health hypertable records periodic health telemetry from each sensing node, indexed by a composite key of timestamp (timestamptz) and node id. Data columns include battery mv (millivolts), rssi dbm (dBm), and uptime s (seconds since last reset). Time-based partitioning keeps range queries efficient as the table accumulates months of deployment data. The acoustic features hypertable stores the preprocessed feature data transmitted by sensing nodes, indexed by timestamp and node id. The primary data column is mfcc (a 15×13 int8-quantized representation of the extracted MFCC vector), accompanied by snr db. Since feature extraction occurs entirely on the sensing node, no raw audio is stored. Records are retained for one year to support longitudinal infestation trend analysis.

\subsubsection{Cluster Layer Tables}

The cluster heads table stores static metadata for each deployed cluster head: cluster id, farm zone, geographic coordinates, install date, and uplink type (Wi-Fi or 4G cellular). It has a one-to-many relationship with both the nodes table and the cluster health hypertable. The cluster health hypertable logs periodic health telemetry from each cluster head, indexed by timestamp and cluster id. It records uptime s, rssi dbm, sd free mb (megabytes of remaining SD storage), and the current uplink type, allowing the dashboard to detect connectivity or storage issues before they cause data loss.

\subsubsection{Cloud Layer Tables}

The classification results hypertable stores the output of the deep learning inference pipeline, indexed by timestamp and node id. Each row contains the predicted label (e.g., ``infested'' or ``healthy''), the confidence score (0.0--1.0), and the model version identifier. Storing the model version with each result supports performance comparisons across model updates. The alert log hypertable records all detection alerts, indexed by timestamp and node id. Each row contains the alert type (e.g., ``infestation detected''), the classification confidence that triggered the alert, and an acknowledged flag indicating whether a plantation manager has reviewed it. Alert logs, along with the nodes and cluster heads metadata tables, are retained indefinitely.

\subsubsection{Data Retention and Partitioning Strategy}

The four hypertables are automatically partitioned by TimescaleDB along the timestamp dimension. Time-range queries, for example, retrieve all classification results for a specific farm zone over the past 30 days, executed against only the relevant partitions rather than the full table. Acoustic features and classification results are retained for one year, after which aged-out partitions are dropped automatically.

\subsection{Detailed Design}

\subsubsection{Dynamic Representation}

The dynamic behavior of the system is described through two operational flowcharts: the Sensing Layer flowchart (Figure 4.5), which details the on-node signal acquisition, processing, and transmission cycle; and the Cluster Layer flowchart (Figure 4.6), which describes the packet reception, validation, buffering, and forwarding logic at each cluster head.

\paragraph{Sensing Layer Operational Flow}

\begin{figure}[H]
\centering
\includegraphics[width=0.88\textwidth,height=0.72\textheight,keepaspectratio]{figures/Sensing_layer_flowchart.pdf}
\caption{Sensing Layer operational flowchart}
\label{fig:SensingLayerFlowchart}
\vspace{0.5\baselineskip}
\end{figure}
\FloatBarrier

Figure \ref{fig:SensingLayerFlowchart} illustrates the complete duty cycle of the sensing node. The recorded signal is passed through a digital bandpass filter (100--2500 Hz) to suppress low-frequency environmental noise such as wind and ground vibrations, as well as high-frequency electronic noise above the band of interest. The filtered signal is squared element-wise to produce a strictly positive power envelope, and a 100 ms moving average is applied to smooth instantaneous noise spikes and reveal the temporal pattern of any larval feeding activity. This processing exploits the characteristic intermittent feeding behavior of RPW larvae, which produces distinct amplitude modulation in the moving average output when larvae are present, in contrast to the relatively flat output from healthy trees \cite{wang2021}. 

An SNR is calculated from the high-level and low-level values of the moving average output. This SNR definition is chosen to avoid false alarms from instantaneous noise spikes, which affect only brief portions of the recording window, while reliably detecting the sustained amplitude modulation characteristic of larval feeding. If the SNR exceeds the 4 dB decision threshold, the node extracts a fixed MFCC representation of 15 frames × 13 coefficients, quantized to int8 with a single global float32 scale factor. The MFCC vector, along with the node identifier and a UTC timestamp, is packaged into a fixed-size LoRa payload of approximately 214 bytes and transmitted in a single packet with an airtime of approximately 440 ms. After transmission, or if the SNR does not exceed the threshold, the node discards the recorded signal and returns to deep sleep. Under this SNR-gated transmission policy, most duty cycles in the absence of infestation produce no transmission, extending battery life and reserving network capacity for genuine detection events.

\paragraph{Cluster Layer Operational Flow}

\begin{figure}[H]
\centering
\includegraphics[width=0.88\textwidth,height=0.72\textheight,keepaspectratio]{figures/Cluster_layer_flowchart.pdf}
\caption{Cluster Layer operational flowchart}
\label{fig:ClusterLayerFlowchart}
\vspace{0.5\baselineskip}
\end{figure}
\FloatBarrier

Figure 4.6 illustrates the operational logic of a single cluster head. On boot, the Raspberry Pi 4 initializes the RFM95W LoRa receiver via SPI, configuring the operating frequency, spreading factor, and bandwidth to match the sensing node transmissions, then enters continuous receive mode.

When a packet arrives, the cluster head checks the header and CRC. Invalid packets, due to transmission corruption, collision, or interference, are discarded immediately. For valid packets, the cluster head parses the payload to extract the node identifier, UTC timestamp, MFCC feature vector, SNR value, and RSSI, then writes the data to local storage on the SD card in an SQLite format as a backup against upstream connectivity failures.

If an uplink connection is available (Wi-Fi or 4G cellular), the data is published to the cloud backend via MQTT as a JSON-encoded message containing the node identifier, timestamp, and MFCC feature vector. If the uplink is temporarily unavailable, the data is buffered locally and retried upon reconnection, implementing the store-and-forward strategy identified in the risk analysis (Section 7.7). After publishing or buffering, the cluster head resumes listening for the next packet.

The cluster head is a stateless relay. It carries out no classification or inference, passes the feature data through unmodified, and operates independently of all other cluster heads. Each cluster head handles its assigned pool of up to 8,000 sensing nodes, so the system scales linearly across plantation zones by deploying additional cluster heads as needed.

\subsection{Classification Model Selection}
The cloud classifier is the final stage of the detection pipeline. Three candidate architectures will be trained on the same dataset and compared: a Convolutional Neural Network (CNN), a CNN-LSTM hybrid, and a Mamba-based selective state space model (SSM).

The CNN and CNN-LSTM architectures are established baselines in RPW acoustic detection and environmental sound classification (Section 2.5). Mamba is a recent selective state-space architecture \cite{gu2024mamba} that has not previously been applied to RPW acoustic data. Evaluating Mamba alongside the established CNN and CNN-LSTM baselines therefore constitutes a novel contribution of this project.

Because a missed infestation (false negative) allows larval damage to continue undetected and undermines the early-detection objective of this project, model selection prioritizes minimizing false negatives. The three candidate models are evaluated using:

Recall - the primary selection metric, quantifying the proportion of genuinely infested trees flagged by the system.

False Negative Rate (FNR) - the complement of recall, reported explicitly to highlight the proportion of missed infestations.

F1 score - the harmonic mean of precision and recall, used as a secondary metric to ensure that high recall is not achieved at the cost of unacceptably high false positive rates.

The candidate that maximizes recall while maintaining an acceptable F1 score will be selected for deployment. The deployed model can be updated over-the-air as new training data becomes available, without requiring physical access to any sensing node.

\section{Standards}

\subsection{Communication Standards}

Cluster heads use the Wi-Fi uplink governed by IEEE 802.11 b/g/n to transfer aggregated MFCC feature packets to a cloud backend. Where Wi-Fi coverage is absent, 4G cellular uplinks use 3GPP LTE standards to transmit data. The physical-layer LoRa transmissions from the sensing nodes to the cluster head conform to Semtech's LoRa modulation specification in operation at the 915 MHz license-free ISM band. The MQTT v3.1.1 protocol specification (ISO/IEC 20922) governs all forwarding of data bound for the cloud. The cloud backend and the dashboard communicate via a RESTful API that complies with HTTP/1.1 (RFC 7230), and TLS 1.2 or later is enforced on all connections (HTTPS, port 443).


\subsection{Data and Encoding Standards}

MQTT and REST API JSON payloads conform to RFC 8259. All timestamps carried in payloads and database records are encoded in ISO 8601 UTC format to ensure consistent time-series alignment across nodes operating within different time zones or DST boundaries.

\subsection{Electrical and Safety Standards}

Signal conditioning blocks using MCP6002 op-amp and the ESP32 based sensing nodes operate at standard low voltage electronics safety margins. Cluster heads powered via mains follow the relevant IEC 60950-1 or IEC 62368-1 compliance for information technology equipment safety standards. Outdoor agricultural deployments adhere to standard low-voltage DC safety practices with a focus on solar-charged sensing node power circuits.

\subsection{Software and Coding Standards}

The ESP32 embedded firmware adopts the Arduino coding conventions as best practices, while PEP 8 is enforced over all Python code in the cloud backend through the Pylint linter. REST APIs are based on a set of architectural constraints (known as REST), whereas JWT authentication is defined according to standards from RFC 7519.

\subsection{Environmental Standards}

Hardware enclosures of all outdoor deployed hardware, whether for sensing nodes or cluster heads, are designed to have at least a minimum ingress protection rating of IP65 (IEC 60529), meaning that they are protected from dust, up to a specified level, and water jets consistent with the heat, humidity and dust conditions found in commercial palm plantations.

\section{Alternative Design Approaches}

The proposed system architecture: low-cost piezoelectric trunk-mounted sensing nodes connected via LoRa to a cluster head that forwards MFCC features to a cloud-based deep learning classifier was selected after evaluating several alternative design approaches identified in the literature review. This section discusses the rationale for the chosen design relative to those alternatives.

\subsection{Invasive Probe-Based Piezoelectric Systems}

Conventional acoustic detection systems use piezoelectric probes inserted through drilled holes in the tree trunk to maximize acoustic coupling between the sensor and the internal wood structure \cite{pinhas2008,karar2022}. Such probe-based systems have demonstrated strong detection performance, including GMM/VQ-based classifiers achieving detection rates of up to 98.9\% in soundproofed laboratory conditions \cite{pinhas2008}. However, this invasive installation approach introduces entry points for secondary pests, pathogens, and fungal infections into each monitored tree, a risk that compounds with the number of sensing nodes deployed across a large plantation. Repeated trunk penetrations also present regulatory and agronomic concerns. The proposed system instead uses the piezoelectric disc in surface-contact mode, placed directly on the bark surface, eliminating the need for trunk penetration and reducing invasiveness while retaining adequate signal sensitivity within the 1.5-2.0 m effective radius documented in the literature \cite{pinhas2008,karar2022}.



\subsection{Non-Acoustic Detection Methodologies}

Three major non-acoustic detection methodologies were evaluated as potential alternatives or complements to acoustic sensing.

\subsubsection{Volatile Organic Compound (VOC) Biosensing}
Volatile Organic Compound (VOC) biosensing exploits infestation-specific metabolic byproducts that diffuse through palm tissue and can be detected externally without destructive sampling. While VOC sensing offers high specificity to larval metabolic activity, current implementations require genetically engineered bacterial biosensor panels and laboratory-grade plate reader equipment that are incompatible with continuous, autonomous field deployment across large plantations. The technology remains at the proof-of-concept stage for RPW detection and does not yet support the real-time, per-tree monitoring resolution required by this project.

\subsubsection{Molecular Marker Detection}
Molecular marker detection using species-specific DNA amplification (PCR) provides highly accurate species confirmation from larval frass or tissue residues recovered from within the trunk. Although molecular methods can detect infestation at very low levels and are not susceptible to acoustic environmental noise, they require periodic manual sample collection, laboratory processing, and reagent handling. This precludes continuous 24/7 autonomous monitoring and makes large-scale deployment labor-intensive and costly, contrary to the automation objectives of this project.

\subsubsection{Pheromone Trapping Systems}
Pheromone trapping systems capture dispersing adult weevils and provide spatial information about infestation hotspots. Although trap optimization studies have demonstrated effective adult capture rates in UAE and Saudi Arabia conditions, pheromone traps are inherently limited to detecting adult insects and yield no information about internal larval activity. Since structural damage to the palm tree occurs during the larval stage, reliance on adult capture alone cannot guarantee detection before irreversible internal damage has occurred, contradicting the early-detection objective of this project.

\subsubsection{Thermal Imaging and Remote Sensing}
Thermal imaging and remote sensing approaches attempt to measure physiological stress signals or internal fermentation heat associated with infestation. However, published accuracy levels for these methods are generally below 75\% in large-scale field assessments, and they are heavily susceptible to environmental noise from ambient temperature variation, making them unsuitable as a primary detection modality.

\subsection{Edge Classification vs. Cloud-Based Classification}

An alternative architectural choice considered was performing the full deep learning classification on the sensing node (edge inference) rather than transmitting features to the cloud. On-device inference using TinyML frameworks would eliminate latency associated with cloud communication and reduce bandwidth requirements to binary infestation/no-infestation flags. However, the CNN/CNN-LSTM hybrid architectures identified as the highest-performing classifiers for RPW detection impose memory and compute requirements that significantly exceed the capabilities of the ESP32 microcontroller (520 KB SRAM, dual-core at 240 MHz) when operating on full MFCC feature vectors. Furthermore, cloud-based inference allows the classification model to be updated over-the-air as new training data becomes available without requiring physical access to deployed sensing nodes, which is a significant operational advantage for large commercial plantations. The proposed design therefore adopts a hybrid strategy: lightweight signal processing, SNR-based gating, and MFCC feature extraction are performed on the ESP32 to drastically reduce the volume of data transmitted. Meanwhile, the computationally intensive deep learning inference is performed on cloud infrastructure with sufficient processing power. This division of labor respects the bandwidth constraints of LoRa and the hardware limits of embedded nodes, while preserving the classification accuracy achievable only with large-scale deep learning architectures.

\subsection{Wi-Fi-Only vs. LoRa-Based Wireless Communication}

A simpler network architecture using only Wi-Fi communication between sensing nodes and a local gateway was considered. Wi-Fi offers high bandwidth and low latency, but its effective range is limited to approximately 50-100 m in open environments, requiring dense gateway placement to cover a large plantation. The power consumption of Wi-Fi radios is also substantially higher than LoRa, making solar or battery-powered field nodes significantly more expensive and complex to engineer. LoRa, by contrast, achieves a range of approximately 1 km per cluster head in plantation conditions while consuming milliwatts during transmission, enabling sparse gateway placement and simple solar-charged power circuits at each sensing node. The lower data rate of LoRa is compatible with the system's design because feature extraction is performed on-node and only compact fixed-size payloads of approximately 214 bytes are transmitted, not raw audio streams. LoRa therefore provides the optimal balance of range, power efficiency, and sufficient throughput for this application.


% ===============================
% CHAPTER 5: Implementation
% ===============================

\chapter{Implementation}
\label{ch:Implementation}

\section{Hardware Implementation}
\label{sec:HardwareImplementation}

Two physical layers are being implemented in the system, a sensing node placed on every tree and a cluster aggregation node per farm zone.

Each sensing node centers around the ESP32-WROOM-32, which handles ADC sampling, signal processing, and LoRa scheduling/transmission. A 27mm piezoelectric disk is mounted on the palm trunk directly to pick up vibrations from larval feeding. The raw output will be in millivolts, so an MCP6002 op-amp will be used to amplify the raw output into a window that the ESP32's ADC can use. A HopeRF RFM95W LoRa module, paired with a 915 MHz RPSMA antenna, handles transmission to the cluster head. Nodes will run on battery with an optional solar charging circuit for longer deployments.

The cluster head is a Raspberry Pi 4 Model B with a second RFM95W for receiving packets. It connects out via Wi-Fi or 4G, runs on mains, and does no signal processing---its only job is to collect packets and forward them upstream.

Circuit schematics, assembly photographs, timing diagrams, and performance metrics (ADC fidelity, op-amp gain, LoRa link budget, power draw) will be documented in Senior Design 2 after bench testing.

\section{Software Implementation}
\label{sec:SoftwareImplementation}


The software design for this project consists of four main classes: Tree, SensingNode, SensorMeasurement, and ClusterHead. Each tree is associated with exactly one sensing node, and each sensing node is associated with exactly one tree. Each SensingNode should have many SensorMeasurement objects, accumulated over time. Each SensingNode is associated with one cluster head, and each cluster head is associated with many nodes (up to 8,000 nodes)

A graphical user interface is planned to be developed, which is intended to be used by farms to view live RPW infestation predictions for each tree, as well as system health data such as node battery life and number of active nodes. 

To communicate between the cluster head and the cloud server,  MQTT broker with TCP/IP will be used. The cluster head should send MFCC vectors to the cloud server along with the most recent system health data. The system health data will be stored in a  system health table in the server hosted PostgresSQL database. A preprocessing pipeline running on the cloud server prepares the MFCC vectors for ML inference. ML inference should run on the cloud server.The ML inference results should be stored in a database hosted on the server. A server script should provide a RESTful API endpoint that allows  the front end of the web Dashboard  to connect with the cloud server., This will allow serving the most recent ML inference results and system health data from the results table and system health table in the PostgresSQL database. Farm employees can login in into the web dashboard through JWT authentication. An email alert system should also be implemented, which uses  SMTP to send automated email alerts to appropriate farm staff  once a predefined confidence threshold for a Tree's infestation is met.

\section{Integration and Deployment}
\label{sec:IntegrationDeployment}

Integration planning is complete, with physical deployment scheduled for Senior Design 2. During that phase, each hardware component will be assembled and validated individually before full system integration.

From the sensing node outward, the ESP32 captures a 10-second acoustic sample, applies bandpass filtering and SNR-based gating, extracts an MFCC feature vector, and transmits a compact payload via LoRa to the cluster head. Communication timing is managed using a TDMA scheme with 900 ms slots. Two transmission windows are defined to match larval activity patterns: nighttime (10 PM--6 AM, hourly) and daytime (6 AM--10 PM, every four hours).

At the cluster head, each valid packet is stored locally on an SD card and then forwarded as a JSON message over MQTT/TCP. In the cloud backend, the received features are processed through the inference pipeline, stored, and evaluated against a confidence threshold. If the threshold is exceeded, an email alert is triggered. A web dashboard provides the plantation manager with a real-time overview of system activity and status.

Physically, each sensing node is installed on a single tree trunk. A single cluster head can cover approximately a 1 km radius, supporting up to 8,000 nodes using Spreading Factor 7. For larger plantations, multiple cluster heads are deployed, each assigned a distinct, non-overlapping range of node IDs, allowing them to operate independently.

Comprehensive integration testing, including multi-node operation, packet loss rates, end-to-end latency from signal capture to alert, and real-world field deployment, will be conducted and reported in Senior Design 2.

\section{User Interface Implementation}
\label{sec:UIImplementation}

The primary user interface for this project is the farm dashboard (prototype screens shown below). The main goal of this dashboard is to present the machine learning inference data in a manner that makes it possible for farms to gain insights on the infestation status of their palm trees. The top of the dashboard screen shows a live summary of the tree measurements, showing the number of trees being monitored, the number of trees suspected to have an infection, and the number of active nodes. Below this summary, the user can see live detailed data on each tree being monitored, with each tree indicated by a card. Further features that could be added to this dashboard are a tree map view, which shows each tree's geolocation on a satellite map, along with the tree's data, allowing farms to easily locate infested trees. Furthermore, each tree card in the main dashboard should have a navigate-to-tree feature, facilitating easy access to the tree's location in the farm.

\FloatBarrier

\begin{figure}[H]
    \centering
    \includegraphics[width=0.9\linewidth]{\detokenize{figures/WhatsApp Image 2026-04-17 at 09.31.17.jpeg}}
    \caption{Farm dashboard prototype - Sensor Map Screen.}
    \label{fig:dashboard_screen_sensormap}
\end{figure}

\begin{figure}[H]
    \centering
    \includegraphics[width=0.9\linewidth]{\detokenize{figures/WhatsApp Image 2026-04-17 at 09.31.17 (1).jpeg}}
    \caption{Farm dashboard prototype - Acoustic Analysis Screen.}
    \label{fig:dashboard_screen_acousticanalysis}
\end{figure}

\begin{figure}[H]
    \centering
    \includegraphics[width=0.9\linewidth]{\detokenize{figures/WhatsApp Image 2026-04-17 at 09.31.18.jpeg}}
    \caption{Farm dashboard prototype screen - Alert Log Screen.}
    \label{fig:dashboard_screen_alertlog}
\end{figure}

\begin{figure}[H]
    \centering
    \includegraphics[width=0.9\linewidth]{\detokenize{figures/WhatsApp Image 2026-04-17 at 09.31.18 (1).jpeg}}
    \caption{Farm dashboard prototype - Historical Records Screen.}
    \label{fig:dashboard_screen_historicalrecords}
\end{figure}

\begin{figure}[H]
    \centering
    \includegraphics[width=0.9\linewidth]{\detokenize{figures/WhatsApp Image 2026-04-17 at 09.31.18 (2).jpeg}}
    \caption{Farm dashboard prototype screen - System Health Screen.}
    \label{fig:dashboard_screen_systemhealth}
\end{figure}

\FloatBarrier

The dashboard prototype screens shown in Figures~\ref{fig:dashboard_screen_sensormap}--\ref{fig:dashboard_screen_systemhealth} were designed using Claude Sonnet 4.6 LLM, using the following prompt (Claude also had previous context about this project outside this prompt):

\begin{quote}
\textit{I am designing a web dashboard for farms to user, it should show stats such as number of trees currently being monitored, trees suspecteted to have infection and the confidence of that infection. It should show the number of nodes currently linked to the system, and the number of measurements we are currently getting. Under these dashboard stats, there should be a view of all trees currently being monitored, with each tree represented as a card, showing infection health of the tree, frequency at which measurements are being done, and the average model infestaton prediciton over last 3 runs. each tree id should be present on the card as well. for accents and ctas, use this green color : \#7CEA9C show  sensor status, historical feature records, classification results, and alert logs to the web-based monitoring dashboard. The single-page dashboard application provides real-time visualization of sensor node and cluster head locations on a map with color-coded infestation status markers, time-series graphs of acoustic energy levels per node, detection event timestamps with confidence scores, and system health indicators including battery levels and connectivity status for all nodes in the network., change the name of that application to be NAKHEL}
\end{quote}


% ===============================
% CHAPTER 6: Validation, Verification, and Performance Analysis
% ===============================

\chapter{Validation, Verification, and Performance Analysis}
\label{ch:Validation}

This section presents the planned approach for evaluating the proposed RPW detection system in Senior I. It describes how the system will be verified against its requirements, how it will be validated for the intended use case, and how its performance will be assessed once implementation is complete.

\section{Testing}
\label{sec:Testing}

Testing will be conducted using both white-box and black-box approaches, depending on the subsystem under evaluation. White-box testing will be used for internal components such as the sensing firmware, signal conditioning logic, feature extraction pipeline, and cloud-side processing modules. Black-box testing will be used for end-to-end system behavior, where input acoustic signals are applied and the resulting classification, transmission, and alert outputs are observed.

Test cases will be derived directly from the system requirements to ensure traceability. Each functional requirement will be mapped to one or more test cases so that the sensing node, wireless communication, classification model, database storage, and dashboard functions are all covered. The planned coverage will include normal operation, noisy acoustic conditions, and simulated infestation-like inputs.

\begin{table}[h]
\centering
\caption{Planned Test Cases by Requirement Area.}
\label{tab:testcases}
\renewcommand{\arraystretch}{1.3}
\begin{tabular}{@{} p{0.25\linewidth} p{0.38\linewidth} p{0.28\linewidth} @{}}
\toprule
\textbf{Requirement Area} & \textbf{Planned Test Focus} & \textbf{Expected Result} \\
\midrule
Sensor acquisition    & Verify acoustic sampling from the trunk-mounted sensor & Signal is captured correctly \\
Signal processing     & Verify filtering and feature extraction                & Noise is reduced and features are generated \\
Wireless transmission & Verify data transfer to the cluster or server          & Packet is received without major loss \\
Classification        & Verify model inference on acoustic input               & Infestation or non-infestation is predicted \\
Alert generation      & Verify notification logic                              & Alert is triggered when threshold is exceeded \\
Dashboard display     & Verify data visualization and status display           & Results appear correctly on the interface \\
\bottomrule
\end{tabular}
\end{table}

\section{Verification and Validation}
\label{sec:VandV}

Verification will confirm that the system is built according to the design specifications. This includes checking that the sensing node samples at the intended rate, that the signal is filtered within the target frequency range, and that the output data is transferred and stored in the correct format. Verification also includes confirming that the database structure supports the required records for sensor readings, classification results, and alerts.

Validation will confirm that the proposed system addresses the actual problem of early RPW detection. The system will be validated by checking whether it can distinguish larval feeding sounds from environmental noise and whether it can support earlier detection than visual inspection. Since this is Senior I, the validation plan focuses on the method to be used rather than final experimental outcomes.

\section{Simulation}
\label{sec:Simulation}

Simulation will be used to evaluate the design before field deployment. Acoustic signals representing healthy trees and infestation-like activity will be used to test the processing pipeline and classification flow. Simulated inputs will also help verify whether alerts are generated only when the confidence level or decision threshold is met.

A simulation-based evaluation will be carried out for the sensing, communication, and software layers. On the hardware side, artificial vibration or recorded signals can be used to mimic larval feeding activity. On the software side, labeled acoustic data can be used to test the classifier, database update logic, and dashboard behavior without requiring live plantation deployment.

\begin{table}[h]
\centering
\caption{Planned Simulation Components.}
\label{tab:simulation}
\renewcommand{\arraystretch}{1.3}
\begin{tabular}{@{} p{0.28\linewidth} p{0.38\linewidth} p{0.25\linewidth} @{}}
\toprule
\textbf{Simulation Component} & \textbf{Input Type} & \textbf{Planned Purpose} \\
\midrule
Sensor node        & Synthetic vibration or recorded acoustic sample & Verify signal capture \\
Filtering stage    & Mixed-frequency audio input                     & Verify noise suppression \\
Feature extraction & Preprocessed acoustic clip                      & Verify feature generation \\
Classifier         & Labeled sample data                             & Verify prediction accuracy \\
Communication layer & Encoded event data                             & Verify packet transmission \\
Dashboard          & Stored detection event                          & Verify monitoring display \\
\bottomrule
\end{tabular}
\end{table}

\section{Performance Metrics}
\label{sec:PerformanceMetrics}

The system will be evaluated using metrics that reflect the objectives of the project. Accuracy will measure how well the classifier distinguishes infestation from non-infestation. Precision and recall will be used to understand false alarms and missed detections. Latency, packet loss, and uptime will be used to assess whether the system is suitable for continuous monitoring in a plantation environment.

\begin{table}[h]
\centering
\caption{Performance Metrics and Their Purpose.}
\label{tab:metrics}
\renewcommand{\arraystretch}{1.3}
\begin{tabular}{@{} p{0.28\linewidth} p{0.65\linewidth} @{}}
\toprule
\textbf{Metric} & \textbf{Purpose} \\
\midrule
Accuracy           & Measures overall classification correctness \\
Precision          & Measures how many positive detections are correct \\
Recall             & Measures how many actual infestation cases are detected \\
False positive rate & Measures incorrect infestation alerts \\
Latency            & Measures time between sensing and alert generation \\
Packet loss        & Measures reliability of wireless communication \\
Uptime             & Measures system stability during continuous operation \\
\bottomrule
\end{tabular}
\end{table}

\section{Performance Analysis Plan}
\label{sec:PerformanceAnalysisPlan}

The performance analysis will compare the behavior of the system across the sensing, communication, and classification stages. The planned analysis will examine whether the acoustic sensor can capture usable signals, whether the wireless network can deliver those signals reliably, and whether the classifier can produce useful detection results in noisy conditions. The final system will also be assessed for responsiveness and monitoring usefulness through the dashboard.

For Senior Design I, the main goal of this section is to define a clear evaluation framework that can be executed in Senior Design II. The project will therefore use this plan to guide implementation, testing, and future experimental analysis.

\chapter{Project Management}
\label{ch:ProjectManagement}

% ===============================
% CHAPTER 7: PROJECT MANAGEMENT
% ===============================

\section{Life Cycle Model}
\label{sec:SWLifecycle}

\subsection{Overview and Rationale}

The development of the IoT-based acoustic detection system for Red Palm Weevil (RPW) infestation presents a multifaceted engineering challenge that integrates hardware prototyping, embedded firmware, wireless networking, cloud infrastructure, and deep learning classification. Given the experimental and iterative nature of this work, particularly the need to continuously refine acoustic sensor configurations and machine learning model performance based on empirical results, the project adopts an Agile-Spiral hybrid software development life cycle model.

This hybrid approach combines the structured risk management and phase-gate discipline of the Spiral Model with the iterative, incremental delivery principles of the Agile methodology. The result is a development process that is both technically rigorous and adaptable, enabling the team to respond to the inherent uncertainties of field-deployable bioacoustic sensing systems.

\subsection{Justification for the Agile-Spiral Model}

Several characteristics of this project make a purely sequential model such as the Waterfall model inappropriate. The detection performance of the acoustic classification system is inherently uncertain at the outset: the quality of signals captured from piezoelectric sensors embedded in palm trunks depends on environmental conditions, sensor placement, and wood properties that can only be evaluated empirically through iterative prototyping and field testing. As documented in Section 1.5, environmental noise from wind, rainfall, and bird activity significantly impacts system performance, and no single pre-defined specification can fully anticipate these effects before testing commences.

Similarly, the deep learning classifiers described in the Literature Review including CNN transfer learning models, BiLSTM networks, and CNN-LSTM hybrid architectures require iterative training, cross-validation, and hyperparameter tuning cycles before reaching the target accuracy of greater than 97\% under the constraints defined in Section 1.4.3. These feedback-driven refinement cycles are fundamentally incompatible with a rigid sequential development process.

At the same time, the project encompasses significant hardware risks, particularly around sensor coupling, wireless transmission reliability, and power management, that benefit from the formal risk identification and mitigation loops characteristic of the Spiral Model. The Agile-Spiral hybrid therefore provides an optimal balance: structured risk management at the macro level, combined with flexible sprint-based development cycles at the implementation level.

\subsection{Model Structure and Phases}

The Agile-Spiral life cycle adopted for this project consists of five major spiral iterations, each structured around four quadrants consistent with Boehm's original Spiral Model formulation: (1) Determine objectives, (2) Identify and resolve risks, (3) Develop and test, and (4) Plan the next iteration. Within the development and test quadrant of each spiral, Agile sprint cycles of two to three weeks are conducted to produce testable increments of the system.

The five major spiral iterations map onto the six major design objectives defined in Section 1.4 of this report, while the last objective remain to ensure operational Integration. This approach serves full traceability between the project’s measurable deliverables and the development process:

\begin{table}[h!]
\centering
\caption{Agile-Spiral Iteration Plan}
\label{tab:spiral_plan}
\renewcommand{\arraystretch}{1.3}
\begin{tabular}{>{\raggedright\arraybackslash}p{2.2cm} >{\raggedright\arraybackslash}p{3.8cm} >{\raggedright\arraybackslash}p{3.8cm} >{\raggedright\arraybackslash}p{1.8cm}}
\toprule
\textbf{Spiral / Phase} & \textbf{Primary Objectives} & \textbf{Key Deliverables} & \textbf{Duration} \\
\midrule
Spiral 1: Requirements \& Architecture &
Finalize system architecture; define sensor node hardware specifications; establish dataset and labeling strategy &
System architecture diagram; hardware BoM; dataset acquisition plan &
3 weeks \\
Spiral 2: Acoustic Sensing Unit &
Prototype piezoelectric sensing node; validate signal quality under lab conditions; implement bandpass filtering pipeline &
Functional sensor node prototype; validated signal acquisition pipeline; lab test report &
4 weeks \\
Spiral 3: Wireless Sensor Network &
Implement WSN architecture; validate LoRa communication; optimize power management; environmental enclosure &
Multi-node WSN deployment; communication reliability report; power consumption analysis &
4 weeks \\
Spiral 4: Deep Learning Classifier &
Train and validate CNN/BiLSTM/hybrid classifiers on TreeVibes dataset; optimize hyperparameters; 10-fold cross-validation &
Trained classification model ($>$97\% accuracy); confusion matrix; cross-validation variance report &
5 weeks \\
Spiral 5: Integration \& Field Validation &
Cloud platform integration; dashboard development; end-to-end system testing; comparative field validation vs. manual inspection &
Deployed IoT system; monitoring dashboard; field validation report; final performance metrics &
4 weeks \\
\bottomrule
\end{tabular}
\end{table}

\FloatBarrier

\subsection{Agile Practices Within Each Spiral Iteration}

Within each spiral iteration, the development team applies core Agile engineering practices to manage task decomposition, progress tracking, and quality assurance. These include:

\begin{itemize}
    \item \textbf{Sprint Planning:} At the start of each two-to-three-week sprint, team members
    identify and assign user stories and engineering tasks relevant to the active spiral
    objective. Tasks are drawn from a maintained product backlog and prioritized based on
    risk criticality and functional dependencies.

    \item \textbf{Iterative Prototyping and Testing:} Hardware and software components are built in
    small, testable increments. Each sprint concludes with a working prototype or testable
    software module. This enables rapid detection of integration failures. For instance,
    incompatibility between the microcontroller sampling rate and the cloud preprocessing
    pipeline before they propagate to later development stages.

    \item \textbf{Continuous Integration and Validation:} The deep learning classification pipeline
    employs automated training and evaluation scripts that are re-run after each model
    architecture modification or dataset augmentation step. Performance metrics including
    accuracy, precision, recall, and F1-score are tracked across sprints to ensure consistent
    progress toward the 97\% accuracy target defined in Section 1.4.3.

    \item \textbf{Sprint Review and Retrospective:} Each sprint concludes with a structured review in
    which completed increments are demonstrated to advisors (Dr. Mohamed Alhajri and Dr.
    Taha Landolsi) for validation against the project objectives. A retrospective session
    identifies process impediments and adjusts task allocation for the subsequent sprint.

    \item \textbf{Risk Register Maintenance:} Consistent with the Spiral Model's risk-driven
    approach, a living risk register is maintained throughout the project. Risks identified
    during each spiral, including environmental noise interference, sensor calibration drift,
    and wireless packet loss under field conditions are formally assessed, prioritized, and
    mitigated before progressing to the next iteration.
\end{itemize}

\subsection{Comparison with Alternative Life Cycle Models}

The Waterfall model was considered but rejected due to its requirement for fully specified, stable requirements prior to development commencement. As the literature review established (Section 2.6), the performance of acoustic classification systems in real-world outdoor environments is heavily dataset- and environment-specific, making it impossible to define precise performance targets and implementation approaches without empirical feedback from prototype testing.

A purely Agile approach such as Scrum was considered as an alternative. While Scrum's sprint-based delivery mechanism is well-suited to the software components of this project, it lacks the explicit risk management framework necessary for managing the hardware uncertainty inherent in deploying piezoelectric sensors in outdoor plantation environments. The Spiral Model's formal risk assessment quadrant addresses this gap directly.

The V-Model was also evaluated but found unsuitable, as it assumes a sequential design process with verification stages tightly coupled to corresponding development stages. The feedback-driven nature of machine learning model development where training, evaluation, and architectural modification cycles are interleaved does not align with the V-Model's linear structure.

The Agile-Spiral hybrid therefore emerges as the most appropriate process model for this project: it provides the iterative, feedback-driven development cycles necessary for machine learning model refinement and hardware prototyping, while imposing the risk discipline and phase-gate structure required to manage a complex, multi-subsystem IoT deployment within a fixed academic timeline.

\subsection{Traceability to Design Objectives}

Each spiral iteration maintains full traceability to the project's design objectives as defined in Section 1.4. The mapping ensures that every development activity is anchored to a measurable, verifiable deliverable:

\begin{itemize}
    \item Spiral 1 supports Objectives 1.4.1 and 1.4.6 by establishing the hardware
    specifications and deployment architecture that govern all subsequent spirals.

    \item Spiral 2 directly addresses Objective 1.4.1 the development and laboratory
    validation of a low-cost acoustic sensing unit capable of detecting larval vibrations within
    a 1.5 to 2.0 meter radius.

    \item Spiral 3 addresses Objective 1.4.2 the implementation of a scalable Wireless Sensor
    Network architecture with validated communication reliability, power management, and
    environmental robustness.

    \item Spiral 4 addresses Objective 1.4.3 the design, training, and validation of a deep
    learning-based acoustic classifier achieving greater than 97\% classification accuracy,
    with 10-fold cross-validation variance below 3\%.

    \item Spiral 5 addresses Objectives 1.4.4 and 1.4.5 the integration of the cloud-based
    monitoring platform and the execution of end-to-end field validation against conventional
    inspection practices.
\end{itemize}

This traceability structure ensures that the life cycle model functions not merely as a development process, but as a quality assurance framework that ties every sprint deliverable back to the project's stated technical and operational requirements.

\section{Work Breakdown Structure}
\label{sec:WBS}

\subsection{Top Level Work}

This Work Breakdown Structure (WBS) decomposes the IoT-Based RPW (Red Palm Weevil) Acoustic Detection System into five major work packages. Each package represents a distinct engineering discipline required to deliver a fully integrated, field-validated solution. The root element the overall detection system branches into: (1) Requirements \& Architecture, which establishes the design blueprint and data strategy; (2) Acoustic Sensing Unit, which covers hardware prototyping and signal conditioning; (3) Wireless Sensor Network, which addresses communication, power, and deployment; (4) Deep Learning Classifier, the core AI pipeline for identifying RPW acoustic signatures; and (5) Integration \& Field Validation, which ties the hardware and software together and proves performance in real-world conditions. The sections below present each work package with its full task breakdown.

\begin{figure}[H]
\centering
\includegraphics[width=0.85\textwidth]{figures/TLP.png}
\caption{Top Level Breakdown}
\label{fig:wbs_toplevel}
\end{figure}

\subsection{WP1 --- Requirements \& Architecture}

This work package establishes the technical and operational foundation for the entire system. It encompasses four tasks: defining the overall system architecture and interface specifications; preparing the hardware bill of materials; planning the acoustic dataset acquisition process; and designing the labeling strategy that will guide supervised model training. Outputs from this package directly inform all subsequent work packages.

\begin{figure}[H]
\centering
\includegraphics[width=0.85\textwidth]{figures/WP1.png}
\caption{WP1 --- Requirements \& Architecture}
\label{fig:wbs_wp1}
\end{figure}

\subsection{WP2 --- Acoustic Sensing Unit}

The Acoustic Sensing Unit translates RPW mechanical vibrations inside palm tree trunks into clean electrical signals ready for wireless transmission. The five tasks span hardware prototyping of piezoelectric sensors, integration of metallic waveguides to amplify acoustic coupling, construction of the signal acquisition pipeline, implementation of bandpass filters to isolate the RPW frequency band, and a controlled lab validation campaign to confirm performance before field deployment.

\begin{figure}[H]
\centering
\includegraphics[width=0.85\textwidth]{figures/WP2.png}
\caption{WP2 --- Acoustic Sensing Unit}
\label{fig:wbs_wp2}
\end{figure}

\subsection{WP3 --- Wireless Sensor Network}

The Wireless Sensor Network (WSN) package ensures that acoustic data collected by individual sensing nodes is reliably transmitted to the cloud processing backend. Tasks address WSN topology design, selection and configuration of LoRa communication modules, power management strategies for extended battery life in remote plantations, weatherproof enclosure design, and the multi-node deployment and integration tests that confirm network-wide reliability.

\begin{figure}[H]
\centering
\includegraphics[width=0.85\textwidth]{figures/WP3.png}
\caption{WP3 --- Wireless Sensor Network}
\label{fig:wbs_wp3}
\end{figure}

\subsection{WP4 --- Deep Learning Classifier}

The Deep Learning Classifier is the intelligence core of the system. It processes raw acoustic data through a multi-stage pipeline: data preprocessing and augmentation to improve model generalization; feature extraction using four complementary spectral representations; and training of three candidate model architectures --- CNN Transfer Learning, BiLSTM, and a CNN-LSTM Hybrid. Hyperparameter tuning and 10-fold cross-validation ensure rigorous, reproducible performance estimation. Task 4.2 (Feature Extraction) is further broken down into four sub-tasks covering MFCC, BFCC, CQCC, and STFT Spectrogram methods.

\begin{figure}[H]
\centering
\includegraphics[width=0.85\textwidth]{figures/WP4.png}
\caption{WP4 --- Deep Learning Classifier}
\label{fig:wbs_wp4}
\end{figure}

\subsection{WP5 --- Integration \& Field Validation}

The final work package closes the loop by deploying the full system end-to-end in real agricultural environments. Seven tasks cover cloud infrastructure provisioning, automated preprocessing pipelines, the model inference and alert engine, a real-time monitoring dashboard for farm operators, field deployment and sensor calibration at palm plantation sites, holistic end-to-end system validation, and a head-to-head comparative analysis against conventional manual inspection methods to quantify the detection advantages offered by the IoT solution.

\begin{figure}[H]
\centering
\includegraphics[width=0.85\textwidth]{figures/WP5.png}
\caption{WP5 --- Integration \& Field Validation}
\label{fig:wbs_wp5}
\end{figure}

\section{Project Schedule and Resource Assignments (Gantt Chart)}
\label{sec:Schedule}

The project schedule for the RPW IoT Detection System is structured across six sequential objectives (O1--O6), spanning the period from January to June 2026. The schedule follows a spiral development methodology, wherein each objective builds upon the outputs of the preceding phase, enabling iterative refinement and early risk mitigation. Table 7.3 presents a detailed breakdown of all deliverables, their corresponding time allocations, and the primary team member responsible for each task, alongside any supporting contributors. The four-member team assumes differentiated roles across hardware, software, networking, and machine learning components, with cross-functional collaboration occurring at key integration milestones. This structured allocation ensures accountability, facilitates progress tracking, and aligns individual workstreams with the overarching system integration timeline.

\renewcommand{\arraystretch}{1.3}
\begin{longtable}{>{\raggedright\arraybackslash}p{6.5cm} >{\raggedright\arraybackslash}p{2.8cm} >{\raggedright\arraybackslash}p{1.8cm} >{\raggedleft\arraybackslash}p{1.8cm}}
\caption{RPW Detection System --- Project Schedule with Resource Assignments (January -- June 2026)}\label{tab:gantt}\\
\toprule
\textbf{Task / Deliverable} & \textbf{Schedule} & \textbf{Lead} & \textbf{Support} \\
\midrule
\multicolumn{4}{l}{\textbf{O1 $\cdot$ Acoustic Sensing Unit (Spiral 1--2)}} \\
Physical Sensor Assembly \& Waveguide & Jan 19 -- Feb 2 & Abdurrahman & --- \\
MCU Programming \& ADC Sampling & Jan 19 -- Feb 2 & Khalid & --- \\
Analog Signal Conditioning \& Pre-filtering & Jan 26 -- Feb 5 & Khalid & --- \\
Bandpass Filter Pipeline 100--2500 Hz & Feb 2 -- Feb 9 & Khalid & --- \\
Outdoor Enclosure \& Rugged Casing & Feb 2 -- Feb 9 & Abdurrahman & --- \\
Lab Signal Validation & Feb 9 -- Feb 16 & Abdurrahman & Khalid \\
\midrule
\multicolumn{4}{l}{\textbf{O2 $\cdot$ Wireless Sensor Network (Spiral 3)}} \\
WSN Topology \& Protocol Design & Feb 16 -- Feb 23 & Kareem & --- \\
LoRa Module Config \& Packet Formatting & Feb 23 -- Mar 5 & Kareem & --- \\
MCU $\leftrightarrow$ Wireless Firmware \& Scheduling & Feb 23 -- Mar 5 & Khalid & Kareem \\
Solar Panel \& Power Regulation Circuit & Feb 23 -- Mar 2 & Abdurrahman & --- \\
Network Reliability \& Sync Testing & Mar 9 -- Mar 16 & Kareem & Abdurrahman, Khalid \\
\midrule
\multicolumn{4}{l}{\textbf{O3 $\cdot$ Deep Learning Classifier (Spiral 4)}} \\
Dataset Preprocessing \& Augmentation & Mar 16 -- Mar 23 & Khalid & --- \\
Feature Extraction --- MFCC / BFCC / CQCC / STFT & Mar 16 -- Mar 26 & Khalid & --- \\
CNN Transfer Learning Model & Mar 23 -- Apr 2 & Haitham & --- \\
BiLSTM Model Architecture & Mar 23 -- Apr 2 & Haitham & --- \\
CNN-LSTM Hybrid Model & Apr 6 -- Apr 13 & Haitham & Khalid \\
Hyperparameter Tuning & Apr 13 -- Apr 20 & Haitham & Khalid \\
10-Fold Cross-Validation (variance $<$ 3\%) & Apr 13 -- Apr 20 & Haitham & Khalid \\
\midrule
\multicolumn{4}{l}{\textbf{O4 $\cdot$ Cloud Platform \& Dashboard (Spiral 5)}} \\
Cloud Server Setup \& Data Ingestion API & Apr 20 -- Apr 27 & Kareem & --- \\
Database Schema \& Security Setup & Apr 20 -- Apr 25 & Kareem & --- \\
Model Inference Service Deployment & Apr 27 -- May 4 & Haitham & --- \\
Automated Preprocessing on Server & Apr 27 -- May 4 & Haitham & --- \\
Alert Threshold Logic \& Delivery & May 4 -- May 9 & Kareem & Haitham \\
Monitoring Dashboard UI & May 4 -- May 11 & Kareem & --- \\
\midrule
\multicolumn{4}{l}{\textbf{O5 $\cdot$ Field Testing \& Validation}} \\
Hardware Durability \& Weatherproofing Checks & May 11 -- May 18 & Abdurrahman & --- \\
Network Stability \& Latency Logging & May 11 -- May 18 & Kareem & --- \\
Signal Quality \& Noise Benchmarking & May 11 -- May 18 & Khalid & --- \\
Model Accuracy vs Manual Inspection & May 11 -- May 18 & Haitham & --- \\
End-to-End System Validation & May 18 -- May 25 & All & --- \\
\midrule
\multicolumn{4}{l}{\textbf{O6 $\cdot$ Field Deployment \& Calibration}} \\
Physical Installation \& Mounting Guide & May 25 -- May 30 & Abdurrahman & --- \\
Network Configuration \& Expansion Guide & May 25 -- May 30 & Kareem & --- \\
Sensor Calibration \& Signal Verification Procedure & May 25 -- May 30 & Khalid & --- \\
Model Retraining \& Update Procedure & May 25 -- May 30 & Haitham & --- \\
End-to-End Integration Guide & May 28 -- Jun 2 & Khalid & Kareem \\
\bottomrule
\multicolumn{4}{l}{\footnotesize Note: ``Lead'' denotes the primary responsible team member; ``Support'' indicates a secondary} \\
\multicolumn{4}{l}{\footnotesize contributor. A dash (---) indicates no supporting member was assigned for that task.} \\
\end{longtable}

\clearpage
\begin{landscape}
\thispagestyle{plain}
\begin{figure}[p]
\centering
\includegraphics[width=\linewidth,height=1.12\textheight,keepaspectratio]{figures/The Gantt Chart.png}
\caption{Project Gantt Chart}
\label{fig:gantt_chart}
\end{figure}
\end{landscape}
\clearpage
\section{Team Organization}
\label{sec:TeamOrg}

The team for this project is composed of three Computer Engineering students, Khalid Hamdan, Kareem Mahmoud, and Abdurrahman Maricar, as well as one Computer Science student, Haitham Ayyash. The group's expertise span a wide range of disciplines, ranging from Wireless communication networking, embedded systems, to Deep Learning.

The group adopts a flat decision making structure,Each group member is expected to find the best solution to the problem they are assigned, and then present their findings in our weekly meetings. Suggestions amongst group members are made during the group's weekly meeting, and a consensus amongst all group members is reached after discussion. The group's joint decision is then consulted with the project's advisors, Dr. Mohamed AlHajri and Dr. Taha Landolsi, for their expert opinion.

Due to the breadth of tasks required by this project, each team member has been assigned specific roles that align with their respective experience. The table below shows a detailed breakdown of the roles of each team member.

For each Objective, the team member can have a lead role, owner role, or support role. The Lead role means that each respective member is primarily responsible for that objective, ensuring that the objective is met to required standards. The Support role means that a team member is responsible for aiding the lead in accomplishing the task. The owner role means the team member is solely responsible for achieving the task, meaning that little to no support from other members is expected.

Note: All team members hold an Owner role for the Field Testing \& Validation and Field Deployment \& Calibration objectives, as each member is individually responsible for distinct tasks within these objectives that fall exclusively under their area of expertise.

\begin{table}[H]
\centering
\caption{Khalid Hamdan --- Roles and Responsibilities}
\label{tab:roles_khalid}
\renewcommand{\arraystretch}{1.3}
\begin{tabular}{>{\raggedright\arraybackslash}p{3.5cm} p{1.8cm} p{8.0cm}}
\toprule
\textbf{Objective} & \textbf{Role} & \textbf{Tasks} \\
\midrule
Acoustic Sensing Unit & Support & Microcontroller programming; ADC sampling rate
configuration; analog signal conditioning circuitry; on-device
pre-filtering pipeline \\
Wireless Sensor Network & Support & MCU wireless module firmware; buffering and transmission
scheduling logic \\
Deep Learning Classifier & Lead & Signal preprocessing pipeline; feature extraction (MFCC,
BFCC, CQCC, STFT spectrograms); dataset preparation and
augmentation; TinyML edge inference feasibility study \\
Field Testing \& Validation & Owner & Signal quality validation in-field; noise condition
benchmarking (wind, birds, machinery) \\
Deployment \& Documentation & Owner & Sensor calibration and signal verification procedure; joint
authorship of system integration guide \\
\bottomrule
\end{tabular}
\end{table}

\begin{table}[H]
\centering
\caption{Kareem Mahmoud --- Roles and Responsibilities}
\label{tab:roles_kareem}
\renewcommand{\arraystretch}{1.3}
\begin{tabular}{>{\raggedright\arraybackslash}p{3.5cm} p{1.8cm} >{\raggedright\arraybackslash}p{8.0cm}}
\toprule
\textbf{Objective} & \textbf{Role} & \textbf{Tasks} \\
\midrule
Wireless Sensor Network & Lead & WSN topology and protocol design; LoRa module
configuration; data packet formatting and transmission;
network reliability and synchronization testing \\
Cloud Platform \& Dashboard & Lead & Cloud server setup and data ingestion API; database schema
for acoustic readings and alerts; network security and access
control; alert delivery (SMS/email/push) \\
Field Testing \& Validation & Owner & Network stability and uptime measurement; data loss rate and
latency logging \\
Deployment \& Documentation & Owner & Network configuration and expansion guide; joint authorship
of system integration guide \\
\bottomrule
\end{tabular}
\end{table}

\begin{table}[H]
\centering
\caption{Abdurrahman Maricar --- Roles and Responsibilities}
\label{tab:roles_abd}
\renewcommand{\arraystretch}{1.3}
\begin{tabular}{>{\raggedright\arraybackslash}p{3.5cm} p{1.8cm} >{\raggedright\arraybackslash}p{8.0cm}}
\toprule
\textbf{Objective} & \textbf{Role} & \textbf{Tasks} \\
\midrule
O1: Acoustic Sensing Unit & Lead & Physical sensor assembly and mounting; piezoelectric
element selection and coupling; metallic waveguide design
and installation; outdoor enclosure and rugged casing \\
O2: Wireless Sensor Network & Support & Physical integration of wireless modules onto node PCB;
solar panel and battery power circuit; power regulation
components \\
O5: Field Testing \& Validation & Owner & Hardware durability and weatherproofing checks; physical
sensor re-mounting and re-calibration \\
O6: Field Deployment \& Calibration & Owner & Physical installation and mounting guide \\
\bottomrule
\end{tabular}
\end{table}

\begin{table}[H]
\centering
\caption{Haitham Ayyash --- Roles and Responsibilities}
\label{tab:roles_haitham}
\renewcommand{\arraystretch}{1.3}
\begin{tabular}{>{\raggedright\arraybackslash}p{3.5cm} p{1.8cm} >{\raggedright\arraybackslash}p{8.0cm}}
\toprule
\textbf{Objective} & \textbf{Role} & \textbf{Tasks} \\
\midrule
O3: Deep Learning Classifier & Lead & CNN/BiLSTM/hybrid model architecture design; model
training, fine-tuning, and hyperparameter search; 10-fold
cross-validation and variance analysis ($<$3\%); confusion
matrix, precision, and recall evaluation \\
O4: Cloud Platform \& Dashboard & Support & Model inference service deployment; automated
preprocessing on server side; alert threshold logic and
confidence scoring \\
O5: Field Testing \& Validation & Owner & Model accuracy comparison against conventional inspection;
false positive and negative analysis on field data \\
O6: Field Deployment \& Calibration & Owner & Model retraining and update procedure \\
\bottomrule
\end{tabular}
\end{table}

\clearpage

\section{Budget Estimation}
\label{sec:Budget}

Table~\ref{tab:bom} presents the bill of materials for one complete sensing node and one cluster head, with all prices in AED. All software tools used in this project are freely available open-source or community editions; no software licensing costs are incurred.

\begin{table}[H]
    \centering
    \caption{Bill of materials and budget estimate (AED).}
    \label{tab:bom}
    \renewcommand{\arraystretch}{1.3}
    \resizebox{\textwidth}{!}{%
    \begin{tabular}{cl>{\raggedright\arraybackslash}p{5.2cm}ccc}
        \toprule
        \textbf{No.} & \textbf{Component} & \textbf{Description} & \textbf{Qty} & \textbf{Unit (AED)} & \textbf{Subtotal (AED)} \\
        \midrule
        \multicolumn{6}{l}{\textit{Sensing Unit — materials required per tree}} \\
		1  & Microcontroller     & ESP32 dev board (Wi-Fi/Bluetooth, dual-core)       & 1 & 25.00 & 25.00 \\
		2  & Piezo sensor        & 27~mm piezoelectric disc                            & 1 & 2.00  & 2.00  \\
		3  & Op-amp IC           & MCP6002 SOP-8                                      & 1 & 1.20  & 1.20  \\
		4  & LoRa module         & HopeRF RFM95W (915 MHz)                     		& 1 & 24.95 & 24.95 \\
		5  & LoRa antenna        & 915~MHz omni antenna (RP-SMA)                      & 1 & 48.67 & 48.67 \\
		6  & Breadboard          & 400-tie breadboard                                 & 1 & 5.00  & 5.00  \\
		7  & Wiring              & Jumper wire pack                                   & 1 & 20.00 & 20.00 \\
		8  & USB cable           & Micro-USB cable                                    & 1 & 10.00 & 10.00 \\
		9 & Battery             & 18650 Li-ion battery (3.7~V)                       & 1 & 2.50  & 2.50  \\
		10 & Voltage regulator   & AMS1117 3.3~V LDO (SOT-223)                        & 1 & 4.00  & 4.00  \\
        \midrule
        \multicolumn{6}{l}{\textit{Cluster Head — one per farm}} \\
		11 & Microcontroller     & Raspberry Pi 4 starter kit (4~GB + 64~GB microSD) & 1 & 560.00 & 560.00 \\
		12 & LoRa module         & HopeRF RFM95W (915 MHz)                    		 & 1 & 24.95  & 24.95  \\
		13 & LoRa antenna        & 915~MHz omni antenna (RP-SMA)                      & 1 & 48.67  & 48.67  \\
		14 & Wiring              & Jumper wire pack                                   & 1 & 7.00   & 7.00   \\
		15 & Storage             & 32~GB microSD card                                 & 1 & 65.00  & 65.00  \\
		16 & Power supply        & 5~V/3~A USB-C wall adapter                         & 1 & 40.00  & 40.00  \\
        \midrule
        \multicolumn{4}{l}{\textit{Software}} & \multicolumn{2}{r}{\textit{All tools free}} \\
        \midrule
		\multicolumn{5}{r}{\textbf{Total}} & \textbf{888.94} \\
        \bottomrule
    \end{tabular}
    }
\end{table}

\section{Engineering Environment and Methodology}
\label{ch:engineering_env}

\subsection{Hardware Environment}
\label{sec:hw_environment}

The hardware development environment covers both physical layers of the system’s dual-layer architecture: (1) sensing node layer and (2) cluster head layer.

\subsubsection{Sensing Node Prototype}

Bench Validation: Each sensing node prototype is assembled on a breadboard (see: Fig. 3A above) and is powered via USB from a laptop. Concerning the processing side of the edge node, it is based on an ESP32 microcontroller and performs ADC sampling, digital signal processing, SNR-based detection gating, MFCC feature extraction at this stage before scheduling a LoRa transmission. Physical coupling of a 27mm piezoelectric disc sensor to a test surface is used to simulate the capture of vibration via trunk mount. Before arriving at the ESP32 ADC, a signal conditioning using an MCP6002 op-amp amplifier circuit is applied to the piezo. The HopeRF RFM95W LoRa transceiver module connects to the ESP32 via SPI and is responsible for all wireless packets sent to the cluster head. The sensing node prototype cost of all components is approximately 143.32 AED.

The bench testing of the sensing node utilizes a signal generator to inject synthetic signals into the amplifier input, and an oscilloscope to confirm ADC sampling fidelity and bandpass filter response before any processing occurs at the software level. End-to-end validation involves verifying that a synthetic larval-band signal triggers the SNR gate followed by an MFCC feature extraction step, and that the cluster head successfully receives a LoRa packet. This yields 12 acoustic analyses per day under the planned monitoring schedule for the prototype.

\subsubsection{Cluster Head Prototype}

The Compact Cluster Head prototype features a Raspberry Pi 4 Model B in conjunction with an RFM95W LoRa module interfaced through SPI. The cluster head is powered by mains all the time (3--4W round-the-clock) so it listens over both transmission windows defined with respect to the sensing layer duty cycle, the nighttime (10PM--6AM) and daytime (6AM--10PM) iterations. The received MFCC feature packets are logged locally onto an SD card and then sent up to the cloud backend over Wi-Fi. The cost of a barebone prototype is approximately 745.62 AED. At the cluster head, no classification is performed since it only serves the purpose of being an always-on aggregation and forwarding tier.

\subsection{Operating System and Runtime Environment}
\label{sec:os_runtime}

Embedded firmware development for ESP32 sensing nodes occurs within the Arduino runtime environment, which offers a hardware abstraction layer compatible with boards based on ESP32. It provides a high level of abstraction that allows for prototyping ADC sampling routines, digital signal processing pipelines and RFM95W LoRa stacks all without the need for low level register configuration.

The Raspberry Pi 4 runs as the cluster head on Raspberry Pi OS (Debian-based, 64-bit). The cluster head receives LoRa packets, logs them to an SD card and forwards the relevant information onto a MQTT broker using Python 3.10 as the primary language in its daemon.

Cloud backend and machine learning inference pipeline are deployed in a Linux-based server environment (ubuntu 22.04 LTS). The data processing pipeline, model training and REST API backend are all written in Python 3.10 Essential Python libraries include librosa for verifying MFCC features, PyTorch for developing the CNN and BiLSTM model, Flask as a REST API server implementation, psycopg2 along with TimescaleDB to interact with our database. Frontend of the web dashboard is developed by standard HTML5, CSS3, and Javascript.

\subsection{Development Tools}
\label{sec:dev_tools}

\subsubsection{Arduino IDE}

The Arduino Integrated Development Environment is used for writing, compiling, and flashing embedded firmware onto the ESP32-based sensing nodes. It provides built-in library management, a serial monitor for real-time debug output, and board manager support for ESP32 targets.

\subsubsection{Visual Studio Code (VS Code)}

For all Python development on the cluster head and cloud backend (signal processing pipeline, machine learning training scripts, REST API backend, web dashboard frontend), the code editor is Visual Studio Code (VS Code). Extensions include a Python extension with Pylance, Jupyter for interactive notebook-based model development; and Pylint to enforce compliance with the PEP8 standard.

\subsubsection{GitHub}

We use GitHub to manage our version control and team collaboration. The source code for everything, ESP32 firmware, cluster head daemon, cloud backend and dashboard, is kept in a private shared repository. Each feature or subsystem is developed on its own branch in a branch-based workflow, but merged into the main branch only after peer review. GitHub Issues will be used to keep track of outstanding tasks and bugs.

\subsection{Development Methodology}
\label{sec:dev_methodology}

The team follows an Agile-inspired iterative development approach which is adjusted for the academic project timeline. At the beginning of every week, there is a well-defined sync meeting where each member shares what has been achieved since the last sync meeting, discusses blockers and decides on tasks to be done that week. This cadence reflects the scrum sprint review and planning cycle without its formal overhead.

Development is structured into four iterative phases that correspond to project milestones: (1) sensing node hardware prototyping and end-to-end LoRa transmission validation to the cluster head, (2) signal processing pipeline implementation including SNR-gated detection and MFCC feature extraction on the ESP32, (3) machine learning model training and evaluation on cloud backend, and finally (4) cluster head forwarding service, cloud infrastructure with dashboard, full end-to-end integration. Deliverables at the end of each phase are tested and validated before moving to the next phase which reduces integration risk and allows for design problems to be detected early.

\section{Risk Analysis}
\label{sec:RiskAnalysis}

This section provides an overview of the potential risks that this project might face, and suggests mitigation strategies to ensure that worst risks are avoided, or ways to address these risks if they occur during the project.

\subsection{Machine Learning Model Accuracy Dropping in Outdoor Field Testing}

This problem has been previously discussed in the literature, where there is concern that machine learning models trained on RPW data collected in indoor laboratory settings may not perform well when deployed in real outdoor farms \cite{pinhas2008}. This is because outdoor environments are subject to many sources of noise, such as wind and birds, that can interfere with the acoustic signals \cite{pinhas2008,rach2013}. Poor field performance would directly undermine the central goal of the project, making mitigation of this risk essential. One mitigation strategy is to collect a dataset of outdoor acoustic measurements and train the model on that data from the outset. While this should help avoid the issue, training on raw noisy data may itself reduce model performance. Therefore, one way to contain this risk is to apply band-pass filtering in the 100--2500 Hz range prior to model inference, ensuring that environmental noise does not significantly hinder classification performance \cite{ashry2022cnn,karar2022,saranya2025}.

\subsection{Training Dataset Not Representative of Early-Stage Infestation}

The TreeVibes dataset \cite{rigakis2021}, which is planned for training, is not explicitly organized around early-stage infestation labels. Since early-stage infestations often do not yet produce visible external symptoms \cite{pinhas2008,sutanto2023}, relying exclusively on TreeVibes may bias the machine learning model toward patterns associated with more advanced infestations, directly undermining a crucial objective of the project.

Therefore, to mitigate this risk, early-stage labeled infestation data need to be obtained and incorporated into model training, either during the initial training stage or through fine-tuning after an initial phase on the TreeVibes dataset. The search for early-stage piezoelectric acoustic data is still underway and remains challenging, as most authors who have demonstrated early-stage detection did so using DAS systems, meaning that their datasets are not compatible with the piezoelectric-based data collected in this project.

If acquiring a small sample of early-stage infestation data is possible, data augmentation techniques can be employed \cite{rezaul2024,turskis} to artificially increase the size of the early-stage dataset. Following this, a transfer learning approach can be adopted, where the model's general acoustic feature learning is first established using the TreeVibes late-stage infestation dataset, and the augmented early-stage data is then used for fine-tuning to improve early detection performance.

\subsection{Losing Wireless Communication Between the Sensor Nodes and the Cluster Head}

Due to the harsh environmental conditions present in palm tree farms, communication between the on-tree sensor nodes and the cluster head may be interrupted, preventing acoustic data from reaching the offsite server for machine learning classification. Losing acoustic data in this way represents a real risk, as it could result in infested trees being entirely missed by the system. To mitigate this risk, a store-and-forward approach should be implemented. First, local storage should be installed on each sensing unit, allowing it to retain recorded data when wireless communication with the cluster head cannot be established. Second, each sensing unit must be capable of automatically reestablishing its wireless connection once communication is restored, at which point the locally stored data can be forwarded to the server. While this mitigation strategy does not directly prevent communication loss from occurring, it ensures that acoustic sensor data is not lost in the case of a communication loss.

\chapter{Conclusions}
\label{ch:conclusions}

\section{Summary}
\label{sec:summary}

This report has presented the design of a low-cost, Internet of Things (IoT)-based acoustic detection system for the early identification of Red Palm Weevil (RPW) (Rhynchophorus ferrugineus) larvae in date palm trees. The system was motivated by the critical gap between infestation onset and visible symptom appearance, which renders conventional visual inspection and pheromone trapping insufficient for early-stage intervention. The proposed solution is intended to address this gap through the integration of non-invasive piezoelectric acoustic sensing, wireless sensor networking, and cloud-based deep learning classification.

The literature review established that acoustic methods are among the most reliable modalities for early RPW detection, with reported classification accuracies exceeding 97\% under controlled laboratory conditions. However, the review also identified key limitations of existing approaches: probe-based piezoelectric systems require invasive trunk penetration and exhibit a restricted sensing radius of 1.5--2.0 metres; Distributed Acoustic Sensing (DAS) systems achieve high performance but at a cost that renders large-scale plantation deployment impractical; and chemical, molecular, and pheromone-based methods, while highly specific, are incompatible with continuous, autonomous real-time monitoring. The proposed system was positioned to address this unmet need by combining low-cost hardware with intelligent signal classification.

The proposed system architecture follows a three-tier design. At the sensing layer, trunk-mounted 27 mm piezoelectric disc sensors would interface with ESP32 microcontrollers, which would perform on-device signal conditioning, bandpass filtering (100--2500 Hz), SNR-based detection gating, and Mel-Frequency Cepstral Coefficient (MFCC) feature extraction. Only acoustically significant events, those exceeding a 4 dB SNR threshold, would be transmitted to preserve bandwidth and extend battery life. At the cluster layer, a Raspberry Pi 4 Model B would receive MFCC feature packets via a LoRa radio link, buffer them locally, and forward them to the cloud backend over Wi-Fi or 4G cellular using MQTT over TCP/IP. At the cloud layer, a PostgreSQL database extended with TimescaleDB would store incoming feature streams alongside classification results and system health telemetry. A trained deep learning model, selected from three candidate architectures (CNN, CNN-LSTM hybrid, and Mamba SSM), would classify each acoustic event and trigger automated email alerts when infestation confidence exceeds a predefined threshold. A web-based dashboard would provide plantation managers with real-time monitoring of sensor status, detection events, and geographic infestation mapping.

The system specifications were formally defined through nine functional and eight non-functional requirements, covering signal acquisition, on-node processing, wireless transmission, cloud classification, database storage, alert generation, and user interface accessibility. The hardware selection was justified against evaluated alternatives, including invasive probe-based systems, DAS, VOC biosensing, molecular marker detection, pheromone trapping, edge classification, and Wi-Fi-only communication, with each alternative found to be unsuitable on grounds of invasiveness, cost, scalability, or operational complexity. An Agile-Spiral hybrid development lifecycle was adopted to manage the inherent uncertainties of field-deployable bioacoustic sensing systems, structured across five spiral iterations that trace directly to the project’s six design objectives.

Physical implementation, field deployment, and end-to-end performance validation are deferred to Senior Design II, at which point the full hardware assembly, multi-node LoRa network operation, and classification model training on the TreeVibes dataset will be completed and evaluated. A structured performance analysis plan, covering accuracy, precision, recall, false negative rate, latency, packet loss, and system uptime, has been defined to guide that evaluation phase.

\section{Future Work}
\label{sec:futurework}

Senior Design I has established the complete system design, architecture, and development plan. The primary continuation in Senior Design II is the physical realization, integration, and validation of the designed system. The following directions constitute the planned and anticipated future work.

The sensing node prototype, consisting of the piezoelectric disc sensor, MCP6002 op-amp signal conditioning circuit, ESP32 microcontroller, and RFM95W LoRa transceiver, will be fully assembled and bench-tested in Senior Design II. Signal fidelity will be verified using a function generator and oscilloscope, and the complete on-device pipeline from ADC sampling through SNR-gated MFCC transmission will be validated before field deployment.

The LoRa-based wireless sensor network will be deployed with multiple sensing nodes operating concurrently under a single cluster head. Network reliability metrics, including packet loss rate, end-to-end latency from signal capture to cloud receipt, and TDMA scheduling stability, will be measured under laboratory and field conditions to confirm that the architecture meets the scalability and reliability requirements defined in Section 3.

Three candidate classification architectures, a CNN transfer learning model, a CNN-LSTM hybrid, and the Mamba selective state-space model, will be trained and evaluated on the TreeVibes acoustic dataset. Model selection will prioritize recall, as missed infestations directly undermine the early-detection objective of the project. Ten-fold cross-validation will be performed, and variance across folds must remain below 3\% in the primary metric to confirm model robustness. The false negative rate and F1 score will be reported as supporting metrics. Evaluating the Mamba architecture alongside the established CNN and CNN-LSTM baselines would constitute a novel contribution of this work, as the architecture has not previously been applied to RPW acoustic classification.

The cloud backend infrastructure, encompassing the MQTT broker, preprocessing pipeline, TimescaleDB database, deep learning inference service, alert engine, and RESTful API, will be provisioned and deployed. The web-based monitoring dashboard will be developed and tested to confirm that real-time sensor status visualization, detection event logging, and automated email alert delivery operate correctly under simulated and live data streams.

Following laboratory integration, the complete system will be deployed in a controlled plantation environment. Performance will be evaluated against the metrics defined in Section 6.4, and a comparative analysis will be conducted against conventional manual visual inspection to quantify the detection advantage offered by the system. This validation will determine whether the system achieves earlier detection than visual inspection and whether it satisfies the scalability and affordability targets established in the design objectives.

The impact of environmental noise sources, including wind speeds exceeding 9 mph, rainfall, bird activity, and agricultural machinery, on classification performance will be evaluated systematically. Strategies for improving robustness, including adaptive SNR thresholding and outdoor dataset augmentation, will be investigated and applied as necessary to ensure reliable operation under real plantation conditions.

Standardized installation protocols, sensor calibration procedures, network configuration guides, and operational manuals will be produced to support the deployment and maintenance of the system by plantation staff without requiring specialist technical expertise.

\printbibliography[title=References]

\end{document}